% Options for packages loaded elsewhere
\PassOptionsToPackage{unicode}{hyperref}
\PassOptionsToPackage{hyphens}{url}
\documentclass[
  12pt,
]{report}
\usepackage{xcolor}
\usepackage[margin=1in]{geometry}
\usepackage{amsmath,amssymb}
\setcounter{secnumdepth}{-\maxdimen} % remove section numbering
\usepackage{iftex}
\ifPDFTeX
  \usepackage[T1]{fontenc}
  \usepackage[utf8]{inputenc}
  \usepackage{textcomp} % provide euro and other symbols
\else % if luatex or xetex
  \usepackage{unicode-math} % this also loads fontspec
  \defaultfontfeatures{Scale=MatchLowercase}
  \defaultfontfeatures[\rmfamily]{Ligatures=TeX,Scale=1}
\fi
\usepackage{lmodern}
\ifPDFTeX\else
  % xetex/luatex font selection
\fi
% Use upquote if available, for straight quotes in verbatim environments
\IfFileExists{upquote.sty}{\usepackage{upquote}}{}
\IfFileExists{microtype.sty}{% use microtype if available
  \usepackage[]{microtype}
  \UseMicrotypeSet[protrusion]{basicmath} % disable protrusion for tt fonts
}{}
\makeatletter
\@ifundefined{KOMAClassName}{% if non-KOMA class
  \IfFileExists{parskip.sty}{%
    \usepackage{parskip}
  }{% else
    \setlength{\parindent}{0pt}
    \setlength{\parskip}{6pt plus 2pt minus 1pt}}
}{% if KOMA class
  \KOMAoptions{parskip=half}}
\makeatother
\usepackage{listings}
\newcommand{\passthrough}[1]{#1}
\lstset{defaultdialect=[5.3]Lua}
\lstset{defaultdialect=[x86masm]Assembler}
\usepackage{longtable,booktabs,array}
\newcounter{none} % for unnumbered tables
\usepackage{calc} % for calculating minipage widths
% Correct order of tables after \paragraph or \subparagraph
\usepackage{etoolbox}
\makeatletter
\patchcmd\longtable{\par}{\if@noskipsec\mbox{}\fi\par}{}{}
\makeatother
% Allow footnotes in longtable head/foot
\IfFileExists{footnotehyper.sty}{\usepackage{footnotehyper}}{\usepackage{footnote}}
\makesavenoteenv{longtable}
\setlength{\emergencystretch}{3em} % prevent overfull lines
\providecommand{\tightlist}{%
  \setlength{\itemsep}{0pt}\setlength{\parskip}{0pt}}
\usepackage{bookmark}
\IfFileExists{xurl.sty}{\usepackage{xurl}}{} % add URL line breaks if available
\urlstyle{same}
\hypersetup{
  hidelinks,
  pdfcreator={LaTeX via pandoc}}

\author{}
\date{}

\begin{document}

\chapter{MScheme: Modula-3/Scheme Interoperability and Comparable
Systems}\label{mscheme-modula-3scheme-interoperability-and-comparable-systems}

Mika Nyström and Claude (Anthropic)\\
March 5, 2026

\section{1. Introduction}\label{introduction}

MScheme is a Scheme interpreter written in Modula-3, shipped with a stub
generator (\passthrough{\lstinline!sstubgen!}) that processes Modula-3
interface files and produces bidirectional bindings. It originated as a
line-for-line translation of Peter Norvig's JScheme (Java) into
Modula-3. The key insight that makes JScheme's Java/Scheme
interoperability seamless is that if the host language has a universal
reference type, then every host-language value already \emph{is} a valid
Scheme value -- no wrapping or boxing is needed. In Java that type is
\passthrough{\lstinline!java.lang.Object!}; in Modula-3 it is
\passthrough{\lstinline!REFANY!}. MScheme carries this principle into
the Modula-3 world.

This report describes MScheme's interoperability architecture in detail,
compares it to other embedding systems that solve similar problems in
different ways, and places the design in historical context -- from its
direct predecessor Tinylisp (DEC SRC, 1989) through the Lisp Machines to
the broader question of how Algol-family languages can support this
architecture on conventional hardware.

\subsection{A taste of MScheme}\label{a-taste-of-mscheme}

Before diving into the architecture, here is real code from
\href{https://github.com/IntelLabs/async-toolkit/tree/main/m3utils/m3utils/csp}{\passthrough{\lstinline!cspc!}},
the CSP compiler built on MScheme. This is Scheme code running inside a
Modula-3 program:

\begin{lstlisting}
;; Read an environment variable -- Env is an M3 module
(define *m3utils* (Env.Get "M3UTILS"))

;; Create M3 BigInt objects -- BigInt.T <: REFANY, so they are
;; ordinary Scheme values
(define a (BigInt.New 12))
(define b (BigInt.New 42))
(BigInt.Equal a b)          ;=> #f
(BigInt.Format a 10)        ;=> "12"

;; M3 procedures used as first-class Scheme values:
;; BigInt.Equal and BigInt.Compare are passed directly to lambda
(define *boolean-binops*
  (list
    `(== integer ,BigInt.Equal)
    `(<  integer ,(lambda (a b) (< (BigInt.Compare a b) 0)))
    `(>  integer ,(lambda (a b) (> (BigInt.Compare a b) 0)))))

;; Format an integer literal by calling M3's Fmt.Int
(define (format-int-literal x)
  (Fmt.Int (BigInt.ToInteger x) 10))

;; Create an M3 object, call its methods from Scheme
(define the-driver
  (obj-method-wrap
    (new-modula-object 'CspCompilerDriver.T)
    'CspCompilerDriver.T))
(the-driver 'init "hello.procs")
(the-driver 'getProcTypes)  ;=> a TextSeq.T (M3 sequence object)
\end{lstlisting}

Every value here -- the \passthrough{\lstinline!BigInt.T!} integers, the
\passthrough{\lstinline!TEXT!} strings, the
\passthrough{\lstinline!CspCompilerDriver.T!} object -- is a native
Modula-3 heap object. None of them are wrapped, copied, or marshaled.
They are Scheme values because every M3 reference type already satisfies
\passthrough{\lstinline!SchemeObject.T = REFANY!}. The procedure names
(\passthrough{\lstinline!BigInt.Equal!},
\passthrough{\lstinline!Fmt.Int!}, etc.) were generated automatically by
\passthrough{\lstinline!sstubgen!} from the corresponding Modula-3
interface files at build time. The generated stub files live in the
platform-specific build directory (e.g.,
\passthrough{\lstinline!csp/ARM64\_DARWIN/FmtSchemeStubs.m3!}).

\begin{center}\rule{0.5\linewidth}{0.5pt}\end{center}

\section{2. MScheme Architecture}\label{mscheme-architecture}

\subsection{2.1 The Universal Value
Type}\label{the-universal-value-type}

The foundation of the entire system is one line of Modula-3:

\begin{lstlisting}
(* SchemeObject.i3 *)
TYPE T = REFANY;
\end{lstlisting}

A Scheme value is \passthrough{\lstinline!REFANY!} -- Modula-3's traced,
polymorphic reference type. Any heap-allocated M3 value (objects, REFs,
TEXTs, arrays) already \emph{is} a valid Scheme object. No wrapper
structs, tag bits on the Scheme side, or special allocation paths. The
M3 runtime's garbage collector manages all values uniformly.

Scheme-specific types are defined as ordinary M3 reference types:

{\def\LTcaptype{none} % do not increment counter
\begin{longtable}[]{@{}
  >{\raggedright\arraybackslash}p{(\linewidth - 2\tabcolsep) * \real{0.1944}}
  >{\raggedright\arraybackslash}p{(\linewidth - 2\tabcolsep) * \real{0.8056}}@{}}
\toprule\noalign{}
\begin{minipage}[b]{\linewidth}\raggedright
Scheme type
\end{minipage} & \begin{minipage}[b]{\linewidth}\raggedright
Modula-3 representation
\end{minipage} \\
\midrule\noalign{}
\endhead
\bottomrule\noalign{}
\endlastfoot
Fixnum & \passthrough{\lstinline!SchemeInt.T = REF INTEGER!} (cached $[-256..256]$) \\
Bignum & \passthrough{\lstinline!Mpz.T!} (GMP arbitrary-precision integer) \\
Exact rational & \passthrough{\lstinline!SchemeRational.T!} (num/den as \passthrough{\lstinline!Mpz.T!}) \\
Inexact real & \passthrough{\lstinline!SchemeLongReal.T = REF LONGREAL!} \\
Arb.-prec.\ float & \passthrough{\lstinline!SchemeMpfr.T!} (MPFR wrapper, configurable precision) \\
Complex & \passthrough{\lstinline!SchemeComplex.T!} (re/im as \passthrough{\lstinline!SchemeObject.T!}) \\
Pair &
\passthrough{\lstinline!SchemePair.T = REF RECORD first, rest : REFANY END!} \\
Symbol & \passthrough{\lstinline!SchemeSymbol.T = Atom.T!} \\
Boolean &
\passthrough{\lstinline!SchemeBoolean.T = BRANDED REF BOOLEAN!} \\
String & \passthrough{\lstinline!SchemeString.T = REF ARRAY OF CHAR!} \\
Procedure & \passthrough{\lstinline!SchemeProcedure.T!} (object with
\passthrough{\lstinline!apply!} method) \\
Any M3 ref & Itself -- already \passthrough{\lstinline!REFANY!} \\
\end{longtable}
}

Runtime type discrimination uses Modula-3's
\passthrough{\lstinline!TYPECASE!} (a type-safe downcast),
\passthrough{\lstinline!ISTYPE!}, and \passthrough{\lstinline!TYPECODE!}
-- the same mechanisms the M3 runtime provides for all M3 code. No
parallel type system is needed.

Modula-3 distinguishes \emph{reference types} (heap-allocated, traced by
the GC) from \emph{value types} (unboxed scalars, records, arrays).
Reference types pass through the Scheme boundary unchanged, but value
types must be boxed into reference types. The
\passthrough{\lstinline!sstubgen!}-generated stubs handle this
automatically:

{\def\LTcaptype{none} % do not increment counter
\begin{longtable}[]{@{}
  >{\raggedright\arraybackslash}p{(\linewidth - 4\tabcolsep) * \real{0.2048}}
  >{\raggedright\arraybackslash}p{(\linewidth - 4\tabcolsep) * \real{0.3614}}
  >{\raggedright\arraybackslash}p{(\linewidth - 4\tabcolsep) * \real{0.4337}}@{}}
\toprule\noalign{}
\begin{minipage}[b]{\linewidth}\raggedright
M3 value type
\end{minipage} & \begin{minipage}[b]{\linewidth}\raggedright
Scheme representation
\end{minipage} & \begin{minipage}[b]{\linewidth}\raggedright
Conversion
\end{minipage} \\
\midrule\noalign{}
\endhead
\bottomrule\noalign{}
\endlastfoot
INTEGER & \passthrough{\lstinline!SchemeInt.T!} &
\passthrough{\lstinline!SchemeInt.FromI!} / \passthrough{\lstinline!ToModula\_INTEGER!} with range check \\
CARDINAL & \passthrough{\lstinline!SchemeInt.T!} &
\passthrough{\lstinline!SchemeInt.FromI!} / \passthrough{\lstinline!ToModula\_CARDINAL!} with bounds check \\
LONGREAL & \passthrough{\lstinline!SchemeLongReal.T!} &
\passthrough{\lstinline!FromLR!} / \passthrough{\lstinline!FromO!} \\
REAL & \passthrough{\lstinline!SchemeLongReal.T!} &
promoted via \passthrough{\lstinline!FLOAT(x, LONGREAL)!} \\
EXTENDED & \passthrough{\lstinline!SchemeLongReal.T!} &
promoted via \passthrough{\lstinline!FLOAT(x, LONGREAL)!} \\
\passthrough{\lstinline!Mpz.T!} & \passthrough{\lstinline!SchemeInt.T!} or \passthrough{\lstinline!Mpz.T!} &
\passthrough{\lstinline!SchemeInt.MpzToScheme!} / \passthrough{\lstinline!ToModula\_Mpz\_T!} \\
\passthrough{\lstinline!Mpfr.T!} & \passthrough{\lstinline!SchemeMpfr.T!} &
\passthrough{\lstinline!SchemeMpfr.FromMpfr!} / \passthrough{\lstinline!ToModula\_Mpfr\_T!} \\
BOOLEAN & \passthrough{\lstinline!SchemeBoolean.T!} &
\passthrough{\lstinline!Truth!} / \passthrough{\lstinline!TruthO!} \\
CHAR & \passthrough{\lstinline!SchemeChar.T!}
(\passthrough{\lstinline!REF CHAR!}) & Cached; all 256 values
preallocated \\
TEXT & \passthrough{\lstinline!SchemeString.T!} &
\passthrough{\lstinline!FromText!} / \passthrough{\lstinline!ToText!} \\
Enumeration & \passthrough{\lstinline!SchemeInt.T!} & As ordinal
index (0, 1, 2, \ldots) \\
Subrange & \passthrough{\lstinline!SchemeInt.T!} & As base type,
with bounds check \\
\end{longtable}
}

Numeric M3 value types map into MScheme's full numeric tower:
exact integers (\passthrough{\lstinline!SchemeInt.T!} for fixnums,
\passthrough{\lstinline!Mpz.T!} for bignums), exact rationals
(\passthrough{\lstinline!SchemeRational.T!}), inexact reals
(\passthrough{\lstinline!SchemeLongReal.T!}), arbitrary-precision
floats (\passthrough{\lstinline!SchemeMpfr.T!}), and complex numbers
(\passthrough{\lstinline!SchemeComplex.T!}).
Small integers $[-256..256]$ and common floats (0.0, 1.0) are cached
to reduce allocation. The \passthrough{\lstinline!SchemeChar!} module
preallocates all 256 character values at initialization.
The \passthrough{\lstinline!SchemeModula3Types!} module provides
the full set of \passthrough{\lstinline!ToModula!}/\passthrough{\lstinline!ToScheme!}
converters, including \passthrough{\lstinline!Mpz.T!} and
\passthrough{\lstinline!Mpfr.T!} as first-class interop types.

Records and arrays require more elaborate treatment because they are
\emph{aggregate} value types. In Modula-3, a
\passthrough{\lstinline!RECORD!} is a stack-allocated or field-embedded
aggregate (like a C struct) and an \passthrough{\lstinline!ARRAY!} is a
contiguous block of elements with fixed bounds. Neither is a reference
-- neither can be \passthrough{\lstinline!REFANY!}. Rather than boxing
them into opaque \passthrough{\lstinline!REF!} wrappers (which would
make the contents inaccessible to Scheme),
\passthrough{\lstinline!sstubgen!} \emph{decomposes} them into Scheme
data structures:

\textbf{Records become association lists.} Each field becomes a
\passthrough{\lstinline!(name . value)!} pair, with the field value
converted recursively. A record
\passthrough{\lstinline!RECORD x: INTEGER; y: REAL END!} with values
\passthrough{\lstinline!x=3, y=1.5!} becomes
\passthrough{\lstinline!((x . 3.0) (y . 1.5))!}. The generated
\passthrough{\lstinline!ToScheme!} converter iterates over fields; the
\passthrough{\lstinline!ToModula!} converter walks the alist, matching
field names by symbol identity and assigning back into a stack-allocated
M3 record:

\begin{lstlisting}
(* Generated by sstubgen for RECORD x: INTEGER; y: REAL END *)
PROCEDURE ToModula_MyRecord(x: SchemeObject.T): MyRecord =
  VAR res: MyRecord; p := SchemePair.Pair(x);
  BEGIN
    WHILE p # NIL DO
      IF    sym_x = First(p.first) THEN res.x := Int(Rest(p.first))
      ELSIF sym_y = First(p.first) THEN res.y := Float(Rest(p.first))
      END;
      p := SchemePair.Pair(p.rest)
    END;
    RETURN res
  END ToModula_MyRecord;
\end{lstlisting}

\textbf{Fixed arrays become lists of (index, element) pairs.} An
\passthrough{\lstinline!ARRAY [0..2] OF INTEGER!} with values
\passthrough{\lstinline!\{10, 20, 30\}!} becomes
\passthrough{\lstinline!((0 . 10.0) (1 . 20.0) (2 . 30.0))!}. The
\passthrough{\lstinline!ToModula!} converter also accepts a simplified
form without explicit indices -- a plain list
\passthrough{\lstinline!(10.0 20.0 30.0)!} -- assigning elements to
successive indices from \passthrough{\lstinline!FIRST!} to
\passthrough{\lstinline!LAST!}. The generated converter handles both
forms:

\begin{lstlisting}
(* Generated by sstubgen for ARRAY [0..2] OF INTEGER *)
PROCEDURE ToModula_A(x: SchemeObject.T): ARRAY [0..2] OF INTEGER =
  VAR res: ARRAY [0..2] OF INTEGER; p := SchemePair.Pair(x);
  BEGIN
    IF p # NIL AND NOT ISTYPE(p.first, SchemePair.T) THEN
      (* Simple form: (10 20 30) *)
      FOR i := FIRST(res) TO LAST(res) DO
        res[i] := Int(p.first); p := SchemePair.Pair(p.rest)
      END
    ELSE
      (* Indexed form: ((0 . 10) (1 . 20) (2 . 30)) *)
      WHILE p # NIL DO
        WITH desc = SchemePair.Pair(p.first) DO
          res[Int(desc.first)] := Int(desc.rest)
        END;
        p := SchemePair.Pair(p.rest)
      END
    END;
    RETURN res
  END ToModula_A;
\end{lstlisting}

\textbf{Multidimensional arrays} are nested. Modula-3's
\passthrough{\lstinline!ARRAY [0..2], [0..3] OF REAL!} is equivalent to
\passthrough{\lstinline!ARRAY [0..2] OF ARRAY [0..3] OF REAL!} -- an
array of arrays. \passthrough{\lstinline!sstubgen!} generates converters
recursively: the outer array's element converter is itself the converter
for the inner array. In Scheme, this becomes a list of lists:
\passthrough{\lstinline!((0 . ((0 . 1.0) (1 . 2.0) ...)) (1 . ((0 . 3.0) ...)) ...)!}.

\textbf{Open arrays} (\passthrough{\lstinline!ARRAY OF T!}, with no
fixed bounds) cannot exist as bare value types -- they exist only as
\passthrough{\lstinline!REF ARRAY OF T!} or as formal parameters. As
\passthrough{\lstinline!REF!} types, they are already
\passthrough{\lstinline!REFANY!} and can be passed through the Scheme
boundary directly; \passthrough{\lstinline!sstubgen!} generates
converters that treat them as Scheme vectors or lists depending on
context.

The key design choice is that \passthrough{\lstinline!sstubgen!}
\emph{decomposes} value types into Scheme-native structures rather than
wrapping them in opaque references. This means Scheme code can inspect
and construct M3 records and arrays using ordinary Scheme operations
(\passthrough{\lstinline!car!}, \passthrough{\lstinline!cdr!},
\passthrough{\lstinline!assq!}), at the cost of a copy at each boundary
crossing.

\subsection{2.2 How M3 Values Enter
Scheme}\label{how-m3-values-enter-scheme}

Because \passthrough{\lstinline!SchemeObject.T = REFANY!}, any M3
reference value can be passed into Scheme simply by assigning it to a
Scheme variable. For a host program embedding MScheme:

\begin{lstlisting}
(* From the host M3 program: *)
VAR interp := NEW(SchemeM3.T).init(paths^);
interp.bind("my-driver", myM3Object);
\end{lstlisting}

The \passthrough{\lstinline!bind!} method defines a symbol in the global
Scheme environment. The M3 object \passthrough{\lstinline!myM3Object!}
becomes directly accessible in Scheme as the value of
\passthrough{\lstinline!my-driver!}. No conversion, copying, or
indirection. Scheme code manipulating this value is manipulating the
exact same M3 heap object.

\subsection{2.3 How Scheme Calls M3
Methods}\label{how-scheme-calls-m3-methods}

MScheme uses a three-layer dispatch mechanism:

\textbf{Layer 1: Type operations registry}
(\passthrough{\lstinline!SchemeProcedureStubs.m3!}, a hand-written
module that provides the registration framework; the individual dispatch
procedures and call stubs plugged into it are generated by
\passthrough{\lstinline!sstubgen!}). For each M3 type that should have
methods callable from Scheme, operations are registered by typecode:

\begin{lstlisting}
SchemeProcedureStubs.RegisterOp(
    TYPECODE(TextSeq.T),    (* M3 runtime typecode *)
    "call-method",           (* operation name *)
    TextSeqCallMethod);      (* M3 procedure to dispatch *)
\end{lstlisting}

The dispatch procedure receives the interpreter, the M3 object, a method
name (as an Atom), and the Scheme argument list. The
\passthrough{\lstinline!interp!} parameter is part of the standard
signature for all dispatch procedures and call stubs; it is passed
through to the individual method stubs, which may need it for exception
handling or Scheme callbacks, even when the dispatcher itself does not
use it directly:

\begin{lstlisting}
PROCEDURE TextSeqCallMethod(interp : Scheme.T;
                            object : SchemeObject.T;
                            args   : SchemeObject.T) : SchemeObject.T
  RAISES { Scheme.E } =
VAR name := Scheme.SymbolCheck(SchemeUtils.First(args));
    rest := SchemeUtils.Rest(args);
BEGIN
  IF name = addhi THEN
    NARROW(object, TextSeq.T).addhi(SchemeUtils.Str(SchemeUtils.First(rest)));
    RETURN SchemeBoolean.True()
  ELSIF name = size THEN
    RETURN SchemeInt.FromI(NARROW(object, TextSeq.T).size())
  ELSIF name = get THEN
    RETURN NARROW(object, TextSeq.T).get(SchemeModula3Types.ToModula_INTEGER(SchemeUtils.First(rest)))
  (* ... *)
  END
END TextSeqCallMethod;
\end{lstlisting}

\textbf{Layer 2: Scheme-level wrapper}
(\passthrough{\lstinline!scheme-lib/src/m3.scm!}). The
\passthrough{\lstinline!obj-method-wrap!} function creates a closure
that dispatches method calls through the type operation registry:

\begin{lstlisting}
(define (obj-method-wrap obj type)
  (lambda args
    (if (eq? (car args) '*m3*)
        obj    ; return the raw M3 object
        (modula-type-op type 'call-method obj (car args) (cdr args)))))
\end{lstlisting}

\textbf{Layer 3: Application code}. The result is natural Scheme syntax:

\begin{lstlisting}
(define seq (obj-method-wrap (new-modula-object 'TextSeq.T) 'TextSeq.T))
(seq 'addhi "hello")
(seq 'size)           ;=> 1
(seq 'get 0)          ;=> "hello"
(seq '*m3*)           ;=> the raw TextSeq.T for passing back to M3
\end{lstlisting}

The \passthrough{\lstinline!'*m3*!} convention lets Scheme code extract
the underlying M3 object to pass it back into M3 methods that expect
typed arguments.

\subsection{2.4 How Scheme Values Return to
M3}\label{how-scheme-values-return-to-m3}

Because Scheme values \emph{are} M3 \passthrough{\lstinline!REFANY!},
return is trivial. An M3 method receiving a Scheme value simply narrows
it:

\begin{lstlisting}
VAR result := interp.evalInGlobalEnv(someExpr);
(* result is REFANY -- narrow to expected type: *)
VAR text := NARROW(result, TEXT);
\end{lstlisting}

If the Scheme code returned a wrapped M3 object, the narrow succeeds and
the caller gets back the original M3 object with its original type. If
it returned a Scheme-specific type (pair, number), the narrow can
distinguish that. \passthrough{\lstinline!ISTYPE!} and
\passthrough{\lstinline!TYPECASE!} handle all cases safely.

\subsection{2.5 The Stub Generator
(sstubgen)}\label{the-stub-generator-sstubgen}

MScheme includes a \emph{stub generator} that processes M3 interface
files at build time and auto-generates bidirectional binding code. This
is the system's most distinctive feature.

\passthrough{\lstinline!sstubgen!} shares substantial code with the
Network Objects stub generator
(\passthrough{\lstinline!m3-comm/stubgen!}). Both are
\passthrough{\lstinline!M3ToolFrame!} extensions that walk the M3 AST
using the M3 Toolkit (m3tk), extract type information via an
\passthrough{\lstinline!AstToType!} module, and generate code from the
type signatures. The network objects \passthrough{\lstinline!stubgen!}
generates marshalling and unmarshalling procedures for remote method
invocation; \passthrough{\lstinline!sstubgen!} adds
\passthrough{\lstinline!TypeTranslator!} and
\passthrough{\lstinline!ValueTranslator!} modules that generate
\passthrough{\lstinline!ToScheme!}/\passthrough{\lstinline!ToModula!}
converters and \passthrough{\lstinline!CallStub!} wrappers for Scheme
interoperability. (See section 8 for the provenance of
\passthrough{\lstinline!sstubgen!} and m3tk.)

The design goal, as stated in the \passthrough{\lstinline!sstubgen!}
documentation, is:

\begin{quote}
``We want to be able to \emph{write Modula-3 code in Scheme.} That is,
anything that can be coded in Modula-3 one should be able to code in
Scheme.''
\end{quote}

The stub generator operates in three phases:

\textbf{Phase 1} parses each M3 interface file using the M3 Toolkit
(m3tk) AST infrastructure. For each procedure, type, constant, and
variable in the interface, it extracts type signatures and generates an
M3 stub module containing:

\begin{itemize}
\tightlist
\item
  \emph{ToModula} converters: functions that take a
  \passthrough{\lstinline!SchemeObject.T!} and produce the appropriate
  M3 type (with type-checking and bounds validation).
\item
  \emph{ToScheme} converters: functions that take an M3 value and
  produce a \passthrough{\lstinline!SchemeObject.T!}.
\item
  \emph{CallStub} procedures: wrappers that unpack Scheme arguments,
  call the real M3 procedure, convert the result, and handle exceptions.
\item
  Registration code that plugs these stubs into the runtime registry.
\end{itemize}

The generated code has two kinds of wrappers: \emph{call stubs} for
interface procedures and \emph{dispatch procedures} (with associated
\emph{method stubs}) for object method calls.

A \textbf{call stub} wraps a single M3 procedure. It receives an
exception handler (\passthrough{\lstinline!excH!}) as an explicit
parameter because the caller (the Scheme evaluator) provides it
directly. The \passthrough{\lstinline!interp!} parameter is needed to
invoke the exception handler via \passthrough{\lstinline!SchemeApply!}
if the M3 procedure raises. A generated call stub looks like:

\begin{lstlisting}
PROCEDURE CallStub_Text_Cat(interp : Scheme.T;
                            args   : SchemeObject.T;
                            excH   : SchemeObject.T) : SchemeObject.T
  RAISES { Scheme.E } =
VAR p := SchemePair.Pair(args);
    a1 := ToModula_TEXT(First(p));   (* Scheme -> TEXT *)
    a2 := ToModula_TEXT(Second(p));
BEGIN
  RETURN ToScheme_TEXT(Text.Cat(a1, a2))  (* TEXT -> Scheme *)
END CallStub_Text_Cat;
\end{lstlisting}

A \textbf{dispatch procedure} routes method calls by name. It does not
take \passthrough{\lstinline!excH!} as a separate parameter; instead,
the exception handler is packed into the \passthrough{\lstinline!args!}
list (as the third element) by the Scheme-level
\passthrough{\lstinline!obj-method-wrap!} caller, and the dispatcher
extracts it before forwarding to the appropriate method stub. The method
stubs themselves have the same signature as call stubs, including the
\passthrough{\lstinline!excH!} parameter.

\textbf{Phase 2} collects all stub modules within a package into a
single \passthrough{\lstinline!RegisterStubs()!} entry point.

\textbf{Phase 3} links all package-level
\passthrough{\lstinline!RegisterStubs!} into the final executable, so
that a single call at startup activates all bindings.

The m3makefile for the CSP compiler (cspc) shows this in practice:

\begin{lstlisting}
import ("sstubgen")
SchemeStubs ("CspExpression")
SchemeStubs ("Fmt")
SchemeStubs ("Atom")
SchemeStubs ("Word")
SchemeStubs ("M3Ident")
SchemeStubs ("FileWr")
(* ... 30+ more interfaces ... *)
ExportSchemeStubs ("cspc")
\end{lstlisting}

Each \passthrough{\lstinline!SchemeStubs!} invocation processes one M3
interface. The result is that Scheme code in cspc can call
\passthrough{\lstinline!Fmt.Int!},
\passthrough{\lstinline!Atom.FromText!},
\passthrough{\lstinline!FileWr.Open!}, etc., as naturally as M3 code
can.

\subsection{2.6 Call Flow for Interface
Procedures}\label{call-flow-for-interface-procedures}

When Scheme code evaluates
\passthrough{\lstinline!(Wr.PutText wr "hello")!}, the call passes
through four layers before reaching the actual M3 procedure. Tracing the
full path:

\textbf{1. Evaluation.} \passthrough{\lstinline!Scheme.Eval!} in
\passthrough{\lstinline!Scheme.m3!} evaluates the operator position. The
symbol \passthrough{\lstinline!Wr.PutText!} resolves to a
\passthrough{\lstinline!SchemePrimitive.T!} object that was registered
at startup. The two argument expressions are evaluated, producing
\passthrough{\lstinline!SchemeObject.T!} values. For two arguments, the
evaluator calls the optimized \passthrough{\lstinline!apply2!} method
directly.

\textbf{2. Primitive dispatch.}
\passthrough{\lstinline!SchemePrimitive.Apply2!} validates the argument
count and calls the main dispatcher (\passthrough{\lstinline!Prims!}).
Interface procedure stubs are registered with IDs above the range of the
built-in primitive enum \passthrough{\lstinline!P!}, so the dispatcher
recognizes them as extensions and calls through to the
\passthrough{\lstinline!ExtDefiner!}:

\begin{lstlisting}
IF t.id > ORD(LAST(P)) THEN
  RETURN NARROW(t.definer, ExtDefiner).apply(t, interp, args)
\end{lstlisting}

\textbf{3. Stub lookup.} \passthrough{\lstinline!ExtDefiner.apply!}
looks up the procedure ID in a hash table, finds a
\passthrough{\lstinline!PrimRec!} containing the qualified name
(\passthrough{\lstinline!Wr!}, \passthrough{\lstinline!PutText!} as
atoms), and calls
\passthrough{\lstinline!SchemeProcedureStubs.CallStubApply!}. This
function walks a linked list of registered stubs, matching the interface
name and item name by atom identity (pointer comparison), and calls the
matching stub procedure.

\textbf{4. Stub execution.} The call stub -- generated by
\passthrough{\lstinline!sstubgen!} at build time -- unpacks arguments
from the Scheme pair list, converts each from
\passthrough{\lstinline!SchemeObject.T!} to the appropriate M3 type
using the \passthrough{\lstinline!ToModula!} functions, calls the real
M3 procedure, converts the result back via
\passthrough{\lstinline!ToScheme!}, and handles any M3 exceptions. The
\passthrough{\lstinline!interp!} parameter is needed to invoke the
Scheme-level exception handler via \passthrough{\lstinline!SchemeApply!}
if the M3 procedure raises; the \passthrough{\lstinline!excH!} parameter
is the exception handler closure itself:

\begin{lstlisting}
PROCEDURE CallStub_Wr_PutText(interp: Scheme.T;
    args, excH: SchemeObject.T): SchemeObject.T RAISES {Scheme.E} =
  VAR wr  := NARROW(First(args), Wr.T);           (* REFANY -> Wr.T *)
      txt := SchemeString.ToText(Second(args));    (* SchemeString -> TEXT *)
  BEGIN
    TRY Wr.PutText(wr, txt);                      (* actual M3 call *)
        RETURN SchemeBoolean.True()
    EXCEPT Wr.Failure(e) =>
        EVAL SchemeApply.OneArg(interp, excH, ...)
    END
  END CallStub_Wr_PutText;
\end{lstlisting}

(All generated call stubs follow the
\passthrough{\lstinline!CallStub\_<Interface>\_<Procedure>!} naming
convention; method stubs use
\passthrough{\lstinline!MethodStub\_<Type>\_<Method>!}.)

For the \passthrough{\lstinline!wr!} argument, which is already a traced
reference (\passthrough{\lstinline!Wr.T!} is an object type), the
conversion is just a \passthrough{\lstinline!NARROW!} -- a type-checked
pointer cast, no data copying. For the \passthrough{\lstinline!txt!}
argument, \passthrough{\lstinline!ToText!} converts the mutable Scheme
string (\passthrough{\lstinline!REF ARRAY OF CHAR!}) to an immutable M3
\passthrough{\lstinline!TEXT!}. The return value is a Scheme boolean.

The entire path -- from \passthrough{\lstinline!Eval!} to the actual
\passthrough{\lstinline!Wr.PutText!} call -- is native M3 procedure
calls. There is no interpreter loop, bytecode dispatch, or reflection.
The overhead beyond a direct M3 call is: argument list traversal
(walking \passthrough{\lstinline!SchemePair!} cells), type conversion
(one \passthrough{\lstinline!NARROW!} + one
\passthrough{\lstinline!ToText!}), and stub lookup (atom comparison in a
linked list, typically short).

\subsection{2.7 Runtime Type Introspection from
Scheme}\label{runtime-type-introspection-from-scheme}

MScheme exposes M3's runtime type system to Scheme via primitives:

\begin{lstlisting}
(rttype-typecode obj)          ; get an object's runtime typecode
(rttype-supertype tc)          ; get a typecode's supertype
(rttype-issubtype? tc1 tc2)    ; subtype relationship
(rttype-maxtypecode)           ; highest allocated typecode
(rtbrand-getname tc)           ; get the brand string for a type
(modula-type-class tc)         ; classification (object, ref, etc.)
(list-modula-type-ops tc)      ; list registered operations
(list-modula-types)            ; list all registered type names
(typename->typecode 'Foo.T)    ; look up typecode by name
\end{lstlisting}

This allows Scheme code to introspect M3 objects, walk the type
hierarchy, and discover available methods -- something that M3 code
itself can do only in limited ways.

\subsection{2.8 Concurrency Model}\label{concurrency-model}

MScheme is designed for multi-threaded M3 programs. The key design rule:
a single \passthrough{\lstinline!Scheme.T!} interpreter instance is
single-threaded, but multiple interpreters may share a global
environment. A global environment is a
\passthrough{\lstinline!SchemeEnvironment.T!} -- the root of a chain of
lexical scoping frames that holds all top-level bindings (primitives,
stub registrations, and user-defined globals). Environments come in
\passthrough{\lstinline!Safe!} (synchronized) and
\passthrough{\lstinline!Unsafe!} (unsynchronized) variants; sharing a
global environment across threads requires the
\passthrough{\lstinline!Safe!} variant.

\subsection{2.9 C Libraries via M3 as
Glue}\label{c-libraries-via-m3-as-glue}

MScheme can be used even when the goal is simply to call C code from
Scheme. Modula-3 serves as a type-safe glue layer: a thin M3 interface
wraps the C functions, and \passthrough{\lstinline!sstubgen!}
auto-exports that interface to Scheme. The programmer never writes
Scheme-specific binding code.

A concrete example is the \passthrough{\lstinline!Mpz!} module in the
\passthrough{\lstinline!m3utils!} tree, which wraps the GNU Multiple
Precision Arithmetic Library (GMP) for use from the CSP compiler's
Scheme environment. The architecture has three layers:

\textbf{Layer 1: C binding.} \passthrough{\lstinline!MpzP.i3!} declares
GMP functions as \passthrough{\lstinline!EXTERNAL!} procedures with C
calling convention:

\begin{lstlisting}
(* MpzP.i3 -- raw C interface *)
<*EXTERNAL "__gmpz_add"*>
PROCEDURE c_add(f0: MpzPtrT; f1: MpzPtrT; f2: MpzPtrT);

<*EXTERNAL "__gmpz_mul"*>
PROCEDURE c_mul(f0: MpzPtrT; f1: MpzPtrT; f2: MpzPtrT);

<*EXTERNAL "mpz_format_decimal"*>
PROCEDURE format_decimal(z: MpzPtrT): Ctypes.const_char_star;
\end{lstlisting}

A small C file (\passthrough{\lstinline!MpzC.c!}) provides helper
functions for formatting via \passthrough{\lstinline!gmp\_asprintf!}.

\textbf{Layer 2: Type-safe M3 wrapper.} \passthrough{\lstinline!Mpz.i3!}
defines an opaque traced type \passthrough{\lstinline!T <: REFANY!} that
wraps the raw \passthrough{\lstinline!mpz\_t!}. The implementation
(\passthrough{\lstinline!MpzOps.m3!}) extracts the address and calls
through to layer 1:

\begin{lstlisting}
(* Mpz.i3 -- public interface *)
TYPE T <: REFANY;
PROCEDURE New(): T;
PROCEDURE add(f0: T; f1: T; f2: T);
PROCEDURE Format(t: T; base := FormatBase.Decimal): TEXT;
\end{lstlisting}

The concrete type is revealed in \passthrough{\lstinline!MpzRep.i3!}:

\begin{lstlisting}
(* MpzRep.i3 -- concrete revelation *)
REVEAL
  Mpz.T = BRANDED Mpz.Brand OBJECT
    val : Val;
  END;

TYPE
  Val = RECORD
    w : ARRAY [0..2] OF Word.T;
    (* this will be cast to an mpz_t, and it is AT LEAST big enough. *)
  END;
\end{lstlisting}

The \passthrough{\lstinline!val!} field embeds the GMP
\passthrough{\lstinline!mpz\_t!} storage directly in the M3 heap object.
The implementation extracts its address for C calls:

\begin{lstlisting}
(* MpzOps.m3 -- implementation *)
PROCEDURE add(f0: T; f1: T; f2: T) =
  BEGIN P.c_add(ADR(f0.val), ADR(f1.val), ADR(f2.val)) END add;
\end{lstlisting}

Because \passthrough{\lstinline!T <: REFANY!},
\passthrough{\lstinline!Mpz.T!} values are traced by the M3 garbage
collector. The opaque type prevents Scheme (or M3) code from accessing
the raw \passthrough{\lstinline!mpz\_t!} memory directly -- only the
operations declared in \passthrough{\lstinline!Mpz.i3!} are available.

Resource cleanup uses Modula-3's \passthrough{\lstinline!WeakRef!}
facility. When \passthrough{\lstinline!Mpz.New!} creates a new wrapper
object, it registers a weak-reference callback via
\passthrough{\lstinline!WeakRef.FromRef!}:

\begin{lstlisting}
(* Mpz.m3 *)
PROCEDURE New() : T =
  VAR res := NEW(T);
  BEGIN
    P.c_init(LOOPHOLE(ADR(res.val), P.MpzPtrT));
    EVAL WeakRef.FromRef(res, CleanUp);
    set_ui(res, 0);
    RETURN res
  END New;

PROCEDURE CleanUp(<*UNUSED*>READONLY w : WeakRef.T; r : REFANY) =
  BEGIN
    WITH this = NARROW(r, T) DO
      P.c_clear(ADR(this.val))
    END
  END CleanUp;
\end{lstlisting}

When the GC determines that an \passthrough{\lstinline!Mpz.T!} object is
unreachable, it schedules the \passthrough{\lstinline!CleanUp!}
callback, which calls GMP's \passthrough{\lstinline!mpz\_clear!} to
release the C-level memory. No explicit deallocation is needed from
either M3 or Scheme code -- the \passthrough{\lstinline!Mpz.T!} values
can be passed around, stored in lists, and forgotten, just like any
other Scheme value. This is the same \passthrough{\lstinline!WeakRef!}
mechanism that Modula-3 programs use for any C resource: FFTW plans,
file descriptors, database handles.

\textbf{Layer 3: Automatic Scheme export.} Adding
\passthrough{\lstinline!SchemeStubs("Mpz")!} to the m3makefile causes
\passthrough{\lstinline!sstubgen!} to generate Scheme bindings for every
procedure in \passthrough{\lstinline!Mpz.i3!}. From Scheme:

\begin{lstlisting}
(define a (Mpz.New))
(define b (Mpz.NewInt 42))
(define c (Mpz.New))
(Mpz.add c a b)
(Mpz.Format c)    ;=> "42"
\end{lstlisting}

The \passthrough{\lstinline!Mpz.T!} values flowing through Scheme are
the same traced M3 heap objects -- no copying or marshaling. The GC
handles their lifetime. The C library's \passthrough{\lstinline!mpz\_t!}
storage is embedded inside the M3 object's record fields and is
deallocated when the M3 object is collected.

This pattern -- C library, M3 interface,
\passthrough{\lstinline!sstubgen!} -- is general. The same approach
could wrap SQLite, zlib, OpenSSL, or any C library. Modula-3's
\passthrough{\lstinline!EXTERNAL!} pragma provides the C binding; its
type system provides safety; \passthrough{\lstinline!sstubgen!} provides
the Scheme bridge. The result is that MScheme functions as a practical
Scheme-with-C-FFI system, even though its design is not about C
interoperability at all.

The \passthrough{\lstinline!Mpz!} module is itself partly
auto-generated: a Scheme script (\passthrough{\lstinline!make-mpz.scm!})
reads GMP's function prototypes and emits
\passthrough{\lstinline!MpzP.i3!}, \passthrough{\lstinline!Mpz.i3!}, and
\passthrough{\lstinline!MpzOps.m3!} -- over 200 functions. This is a
second level of code generation: Scheme generates M3 code, then
\passthrough{\lstinline!sstubgen!} generates Scheme stubs from that M3
code, closing the loop.

\subsection{2.10 Summary of Key
Properties}\label{summary-of-key-properties}

\begin{itemize}
\tightlist
\item
  \textbf{Zero-copy value sharing}: M3 objects are Scheme objects, no
  wrapping needed.
\item
  \textbf{Unified GC}: Both languages share the M3 traced heap and
  garbage collector. No reference counting, prevent-collection APIs, or
  weak reference coordination.
\item
  \textbf{Automatic stub generation}: Build-time processing of M3
  interfaces produces bidirectional bindings with type checking.
\item
  \textbf{Runtime type introspection}: Scheme code can query the M3 type
  system, discovering types and operations dynamically.
\item
  \textbf{Method dispatch via typecode}: O(n) scan of registered
  operations per typecode, then atom comparison for method name.
\end{itemize}

\begin{center}\rule{0.5\linewidth}{0.5pt}\end{center}

\section{3. Comparable Systems}\label{comparable-systems}

\subsection{3.1 Lua/C (Lua C API)}\label{luac-lua-c-api}

\textbf{Core representation}: A \emph{stack} of tagged union values
(\passthrough{\lstinline!TValue!} internally). The C API never exposes
\passthrough{\lstinline!TValue!} directly; all value exchange is through
positional stack slots.

\textbf{Wrapping C objects}:
\passthrough{\lstinline!lua\_newuserdata(L, size)!} allocates a
GC-managed memory block. A \emph{metatable} (a Lua table) is attached to
define behavior. The \passthrough{\lstinline!\_\_index!} metamethod
routes field/method lookups to a table of C functions.
\passthrough{\lstinline!\_\_gc!} provides a destructor.

\textbf{Method dispatch}: When Lua code calls
\passthrough{\lstinline!obj:method(args)!}, the VM performs a metatable
lookup for \passthrough{\lstinline!\_\_index!}, finds the method in the
method table, and calls the C function with the Lua state. The C
function pops arguments from the stack and pushes results.

\textbf{Values back to C}: Typed accessors read stack slots:
\passthrough{\lstinline!lua\_tonumber!},
\passthrough{\lstinline!lua\_tostring!},
\passthrough{\lstinline!lua\_touserdata!}.

\textbf{Registration}: C functions are registered via arrays of
\passthrough{\lstinline!luaL\_Reg!} structs, or individually with
\passthrough{\lstinline!lua\_pushcfunction!}.

\textbf{Key trade-offs}: The stack-based API provides complete ABI
isolation -- Lua's internal value representation can change without
breaking extensions. But every boundary crossing requires explicit
push/pop operations. There is no automatic type mapping, stub
generation, or introspection of the host's type system.

\subsection{3.2 GNU Guile (Scheme in C)}\label{gnu-guile-scheme-in-c}

\textbf{Core representation}: \passthrough{\lstinline!SCM!}, a
machine-word-sized tagged value. Fixnums, characters, and booleans are
encoded in tag bits. Heap objects (pairs, foreign objects, strings) are
tagged pointers.

\textbf{Wrapping C objects}: \emph{Foreign object types} replace the
older SMOB mechanism.
\passthrough{\lstinline!scm\_make\_foreign\_object\_type!} defines a
type with named slots and a finalizer. Instances store C pointers in
slots.

\textbf{Method dispatch}: Typically via named procedures
(\passthrough{\lstinline!my-object-get-x!}). Optional GOOPS integration
(CLOS-style generic functions with multiple dispatch) is possible but
uncommon for foreign objects.

\textbf{Values back to C}: \passthrough{\lstinline!scm\_to\_int(val)!},
\passthrough{\lstinline!scm\_to\_double(val)!}, etc. Every function
takes and returns \passthrough{\lstinline!SCM!}.

\textbf{Registration}:
\passthrough{\lstinline!scm\_c\_define\_gfunc("name", nreq, nopt, nrest, cfunc)!}
registers a C function as a Scheme procedure in the current module.

\textbf{Key trade-offs}: The uniform \passthrough{\lstinline!SCM!} type
is simpler than Lua's stack (no positional bookkeeping), but every C
function must do its own type checking. Guile's conservative GC (BDW-GC)
is very friendly to C embedders -- no need to track roots manually. But
Guile provides no automatic stub generation from C header files.

\subsection{3.3 Tcl/C}\label{tclc}

\textbf{Core representation}: \passthrough{\lstinline!Tcl\_Obj*!}, a
dual-ported value with both a string representation and a cached native
(``internal'') representation. Conversions between the two forms are
lazy and cached. Reference counted (no GC).

\textbf{Wrapping C objects}: Tcl uses the ``command-as-object'' pattern.
Creating an object registers a new Tcl command named after it. The
\passthrough{\lstinline!ClientData!} (a \passthrough{\lstinline!void*!})
stored with the command points to the C struct. ``Methods'' are
subcommands: \passthrough{\lstinline!.button configure -text "OK"!}.

\textbf{Method dispatch}: String-based. The first argument to the
object's command is the subcommand name, matched by string comparison or
hash lookup. Object identity is a string name.

\textbf{Registration}:
\passthrough{\lstinline!Tcl\_CreateObjCommand(interp, "name", func, clientData, deleteProc)!}.

\textbf{Key trade-offs}: ``Everything is a string'' gives extraordinary
composability but loses type safety. Reference counting provides
deterministic destruction but cannot handle cycles. The
command-as-object pattern is elegant and uniform but means object
identity is a string, and there is no direct pointer from Tcl to C
memory without going through command lookup.

\subsection{3.4 CPython Extension API}\label{cpython-extension-api}

\textbf{Core representation}: \passthrough{\lstinline!PyObject*!}, a
heap-allocated struct with a reference count and a pointer to a
\passthrough{\lstinline!PyTypeObject!}.

\textbf{Wrapping C objects}: The embedder fills out a
\passthrough{\lstinline!PyTypeObject!} struct with dozens of slots:
\passthrough{\lstinline!tp\_methods!} (method table),
\passthrough{\lstinline!tp\_getset!} (computed properties),
\passthrough{\lstinline!tp\_members!} (direct struct field access by
offset), \passthrough{\lstinline!tp\_dealloc!} (destructor), plus slots
for numeric, sequence, mapping, iterator, and async protocols.

\textbf{Method dispatch}: Attribute lookup via the type's method
resolution order (MRO). If the name is found in
\passthrough{\lstinline!tp\_methods!}, a bound method object is created
wrapping the C function.

\textbf{Values back to C}:
\passthrough{\lstinline!PyLong\_AsLong(obj)!},
\passthrough{\lstinline!PyFloat\_AsDouble(obj)!},
\passthrough{\lstinline!PyArg\_ParseTuple(args, "is|d", ...)!}.

\textbf{Registration}: Module initialization fills a
\passthrough{\lstinline!PyModuleDef!} struct and registers types with
\passthrough{\lstinline!PyType\_Ready!}.

\textbf{Key trade-offs}: Maximum control over Python's internals, but
the reference counting discipline is notoriously error-prone. Every
boundary crossing allocates heap objects. Defining even a simple type
requires filling many \passthrough{\lstinline!PyTypeObject!} slots. The
API is primarily designed for \emph{extending} Python, not for
\emph{embedding} it.

\subsection{3.5 Kawa Scheme (JVM)}\label{kawa-scheme-jvm}

\textbf{Core representation}:
\passthrough{\lstinline!java.lang.Object!}. Kawa compiles Scheme to JVM
bytecode; Scheme values are Java objects.

\textbf{Wrapping Java objects}: No wrapping needed. Java objects are
directly usable in Kawa Scheme.
\passthrough{\lstinline!(define pt (java.awt.Point 10 20))!} creates a
Java \passthrough{\lstinline!Point!}. \passthrough{\lstinline!pt:x!}
accesses the field.

\textbf{Method dispatch}: When the type is known at compile time, Kawa
generates \passthrough{\lstinline!invokevirtual!} bytecode -- as
efficient as Java. When unknown, it falls back to reflection. Type
annotations (\passthrough{\lstinline!::java.awt.Point!}) guide the
compiler.

\textbf{Values back to Java}: Scheme values are already JVM objects.
Java code invokes Scheme functions by calling
\passthrough{\lstinline!apply!} on
\passthrough{\lstinline!gnu.mapping.Procedure!} objects.

\textbf{Registration}: None needed -- both languages share the JVM class
loader and object space.

\textbf{Key trade-offs}: Zero-cost interoperability when types are
known, but reflection fallback is slow. Kawa must work within JVM
constraints (no true tail calls or first-class continuations without
trampolining). The shared-runtime model is fundamentally different from
all other systems here.

\subsection{3.6 JScheme (Scheme in Java)}\label{jscheme-scheme-in-java}

\textbf{Core representation}:
\passthrough{\lstinline!java.lang.Object!}. JScheme is a tree-walking
Scheme interpreter written in Java. In Norvig's original (1998), all
numbers are \passthrough{\lstinline!java.lang.Double!}, symbols are
interned \passthrough{\lstinline!java.lang.String!}, strings are
\passthrough{\lstinline!char[]!} (for mutability), pairs are a custom
\passthrough{\lstinline!Pair!} class with
\passthrough{\lstinline!Object first, rest!} fields, and
\passthrough{\lstinline!null!} represents the empty list. The later
Anderson/Hickey version (2003+) switched numbers to
\passthrough{\lstinline!Integer!} or \passthrough{\lstinline!Double!} as
appropriate, symbols to a custom \passthrough{\lstinline!Symbol!} class,
and strings to \passthrough{\lstinline!java.lang.String!}.

\textbf{Wrapping Java objects}: No wrapping needed. Any Java object is
already an \passthrough{\lstinline!Object!} and therefore a valid Scheme
value. The evaluator treats anything that is not a symbol and not a pair
as self-evaluating, so a \passthrough{\lstinline!java.awt.Frame!} or
\passthrough{\lstinline!java.util.HashMap!} placed into a Scheme
variable simply evaluates to itself.

\textbf{Method dispatch}: Entirely reflection-based. The later JScheme
provides a dot-notation syntactic sugar parsed at read time:

\begin{lstlisting}
(import "java.awt.Font")
(Font. "Serif" Font.BOLD$ 24)     ; constructor (trailing dot)
(.setFont component font)         ; instance method (leading dot)
(.name$ obj)                      ; field access (trailing $)
(Math.round 123.456)              ; static method (dot in middle)
\end{lstlisting}

These desugar into calls to \passthrough{\lstinline!JavaMethod!},
\passthrough{\lstinline!JavaConstructor!}, or
\passthrough{\lstinline!JavaField!} objects, all of which use
\passthrough{\lstinline!java.lang.reflect!} internally. Method tables
are cached per class in hash tables, and overload resolution follows
Java's ``most specific applicable method'' rule at runtime.

\textbf{Values back to Java}: Scheme values are already
\passthrough{\lstinline!Object!}. Java code calls
\passthrough{\lstinline!JScheme.call("proc-name", arg1, arg2)!} and
receives an \passthrough{\lstinline!Object!} result, narrowing via cast.

\textbf{Scheme as Java interface}: The
\passthrough{\lstinline!SchemeInvocationHandler!} class uses
\passthrough{\lstinline!java.lang.reflect.Proxy!} to let a Scheme
closure implement any Java interface, so Scheme callbacks work anywhere
Java expects an interface (listeners,
\passthrough{\lstinline!Runnable!}, etc.).

\textbf{Registration}: \passthrough{\lstinline!(import "java.awt.*")!}
makes classes available; dot-notation symbols are resolved against the
import table. From the Java side,
\passthrough{\lstinline!JScheme.setGlobalValue("name", obj)!} injects
values.

\textbf{Key trade-offs}: JScheme achieves the same zero-wrapping
interoperability as Kawa -- Java objects \emph{are} Scheme values -- but
pays for it with reflection overhead on every method call. Kawa compiles
to \passthrough{\lstinline!invokevirtual!} bytecode; JScheme interprets
and reflects. The simplicity is striking (Norvig's original is
\textasciitilde1700 lines), but the performance gap is large. JScheme is
the direct ancestor of MScheme and establishes the key principle that
MScheme inherits: when the host language has a universal reference type
and a GC, the embedding boundary can be made nearly invisible.

\subsection{3.7 Common Lisp / CFFI}\label{common-lisp-cffi}

\textbf{Core representation}: Tagged Lisp objects
(implementation-dependent tagging scheme). Foreign values are
\passthrough{\lstinline!cffi:foreign-pointer!} wrappers.

\textbf{Wrapping C objects}: The programmer wraps foreign pointers in
CLOS objects and defines generic functions that call through to C.
\passthrough{\lstinline!trivial-garbage:finalize!} attaches destructors.

\textbf{Method dispatch}: Generic functions with CLOS multiple dispatch.
Methods dispatch on the classes of \emph{all} arguments, not just the
receiver.

\textbf{Values back to C}:
\passthrough{\lstinline!cffi:foreign-slot-value!} reads struct fields by
declared offset. \passthrough{\lstinline!defcfun!} defines C function
bindings with automatic marshaling.

\textbf{Registration}: \emph{Pull-based} -- the Lisp side declares C
bindings with \passthrough{\lstinline!defcfun!},
\passthrough{\lstinline!defcstruct!},
\passthrough{\lstinline!defcenum!}. The C library is unmodified.

\textbf{Key trade-offs}: CFFI is portable across CL implementations but
requires manual struct layout declarations and memory management. CLOS
provides a powerful ``glue'' layer but is per-project, not standardized
for FFI.

\subsection{3.8 Racket FFI}\label{racket-ffi}

\textbf{Core representation}: Tagged Racket objects. C types are
first-class Racket values: \passthrough{\lstinline!\_int!},
\passthrough{\lstinline!\_double!}, \passthrough{\lstinline!\_pointer!},
\passthrough{\lstinline!\_fun!}.

\textbf{Wrapping C objects}:
\passthrough{\lstinline!define-cpointer-type!} creates tagged pointer
types. \passthrough{\lstinline!define-cstruct!} auto-generates
constructors, accessors, and mutators. Custom type combinators
(\passthrough{\lstinline!make-ctype!}) attach conversion functions to
base C types.

\textbf{Method dispatch}: Not object-oriented at the FFI layer.
Operations on foreign values are called as procedures.

\textbf{Values back to C}: \passthrough{\lstinline!\_fun!} type
descriptors handle all marshaling declaratively.
\passthrough{\lstinline!(\_fun \_string \_int -> \_bool)!} specifies
argument and return conversions.

\textbf{Registration}: \emph{Pull-based} --
\passthrough{\lstinline!get-ffi-obj!} or
\passthrough{\lstinline!define-ffi-definer!} declares bindings. The C
library is unmodified.

\textbf{Key trade-offs}: The most algebraic and composable FFI design.
Types as values enable declarative marshaling. Tagged pointers add
safety. But the C library must be described entirely from the Racket
side, and per-call marshaling overhead is significant for tight loops.

\subsection{3.9 Tinylisp (Modula-2+)}\label{tinylisp-modula-2}

\textbf{Core representation}: \passthrough{\lstinline!REFANY!}. Tinylisp
(DEC Systems Research Center, 1989) is a Lisp implemented in Modula-2+,
the SRC dialect of Modula-2 that added garbage collection, exception
handling, \passthrough{\lstinline!REFANY!}, and
\passthrough{\lstinline!TYPECASE!}. Symbolic expressions are ordinary
Modula-2+ ref types: \passthrough{\lstinline!Ref.Integer!},
\passthrough{\lstinline!Ref.LongReal!},
\passthrough{\lstinline!Ref.Boolean!},
\passthrough{\lstinline!Ref.Char!}, \passthrough{\lstinline!Text.T!},
\passthrough{\lstinline!SxSymbol.T!}, \passthrough{\lstinline!List.T!},
\passthrough{\lstinline!Ref.Vector!}. Any application's opaque-ref type
is directly usable as a Tinylisp value -- a
\passthrough{\lstinline!Wire.T!} or \passthrough{\lstinline!Wr.T!} is a
first-class s-expression, just as an \passthrough{\lstinline!Mpz.T!} is
a first-class MScheme value.

\textbf{Wrapping Modula-2+ objects}: No wrapping needed. Any traced
reference is already \passthrough{\lstinline!REFANY!} and therefore a
valid Tinylisp value.

\textbf{Method dispatch}: Via generated stubs. A \emph{stub compiler}
(\passthrough{\lstinline!compile\_tli!}) processes hand-written
\passthrough{\lstinline!.tli!} (TinyLisp Interface) files and generates
Modula-2+ stub procedures that dynamically check argument counts and
types, convert between s-expression types and Modula-2+ value types, and
call the corresponding procedures. A \passthrough{\lstinline!.tli!} file
declares modules, types, exceptions, and procedures in s-expression
syntax:

\begin{lstlisting}[language=Lisp]
(MODULE Text Lisp)
(TYPE T)
(EXCEPTION EndOfFile)
(PROCEDURE FromChar (CHAR) "Text.T")
(PROCEDURE Length ("Text.T") INTEGER)
(PROCEDURE SubText ("Text.T" INTEGER & INTEGER) "Text.T")
\end{lstlisting}

\textbf{Values back to Modula-2+}: Scheme values are already
\passthrough{\lstinline!REFANY!}. The stubs narrow them via
\passthrough{\lstinline!TYPECASE!} for value types
(\passthrough{\lstinline!INTEGER!}, \passthrough{\lstinline!CHAR!},
\passthrough{\lstinline!BOOLEAN!}, etc.) and
\passthrough{\lstinline!NARROW!} for reference types.

\textbf{Registration}: The generated stubs register procedures at
startup, similar to MScheme's approach.

\textbf{Execution model}: Unlike MScheme's tree-walking interpreter,
Tinylisp compiled expressions directly into machine code at runtime (an
on-the-fly compiler), making Tinylisp-to-Tinylisp calls roughly as fast
as Modula-2+ calls.

\textbf{Key trade-offs}: Tinylisp achieves the same zero-wrapping,
shared-heap architecture as MScheme -- the underlying mechanism
(\passthrough{\lstinline!REFANY!} + tracing GC +
\passthrough{\lstinline!TYPECASE!}) is identical. The key difference is
in stub generation: Tinylisp requires a hand-written
\passthrough{\lstinline!.tli!} file (essentially a manual
transliteration of the Modula-2+ \passthrough{\lstinline!.def!} file
with type annotations), while MScheme's
\passthrough{\lstinline!sstubgen!} reads the standard M3
\passthrough{\lstinline!.i3!} interface files directly via m3tk.
Tinylisp is also faster at the Lisp level (compiled vs.~interpreted) but
was designed as a lightweight extension language rather than a full
Scheme.

\begin{center}\rule{0.5\linewidth}{0.5pt}\end{center}

\section{4. Comparative Analysis}\label{comparative-analysis}

\subsection{4.1 The Central Trade-off: Shared Heap vs.~Marshaling
Boundary}\label{the-central-trade-off-shared-heap-vs.-marshaling-boundary}

The most fundamental architectural difference across these systems is
whether the host and guest languages \emph{share} their object
representations or maintain \emph{separate} representations connected by
a marshaling boundary.

\textbf{Shared heap} (MScheme, Tinylisp, JScheme, Kawa): The host
language's objects are directly usable as scripting language values. No
per-call conversion cost. The garbage collector sees everything. This is
only possible when both languages are GC'd and share a runtime (M3 +
MScheme, Modula-2+ + Tinylisp) or a VM (Java + JScheme, Java + Kawa).

\textbf{Marshaling boundary} (Lua, Guile, Tcl, Python, CFFI, Racket):
The host creates scripting-language values via conversion APIs. Every
boundary crossing allocates and converts. The two sides maintain
separate type systems and (often) separate memory management.

MScheme's \passthrough{\lstinline!SchemeObject.T = REFANY!} and
Tinylisp's use of \passthrough{\lstinline!REFANY!} are the extreme cases
of shared heap: it is not that host objects are \emph{wrapped} as
Lisp/Scheme values, but that they literally \emph{are} Lisp/Scheme
values, by type identity.

\subsection{4.2 Type Safety and Dispatch
Cost}\label{type-safety-and-dispatch-cost}

{\def\LTcaptype{none} % do not increment counter
\begin{longtable}[]{@{}
  >{\raggedright\arraybackslash}p{(\linewidth - 4\tabcolsep) * \real{0.1163}}
  >{\raggedright\arraybackslash}p{(\linewidth - 4\tabcolsep) * \real{0.4651}}
  >{\raggedright\arraybackslash}p{(\linewidth - 4\tabcolsep) * \real{0.4186}}@{}}
\toprule\noalign{}
\begin{minipage}[b]{\linewidth}\raggedright
System
\end{minipage} & \begin{minipage}[b]{\linewidth}\raggedright
Type check at boundary
\end{minipage} & \begin{minipage}[b]{\linewidth}\raggedright
Method dispatch cost
\end{minipage} \\
\midrule\noalign{}
\endhead
\bottomrule\noalign{}
\endlastfoot
MScheme & \passthrough{\lstinline!TYPECASE!} (M3 runtime typecheck) &
Atom comparison per registered op \\
Tinylisp & \passthrough{\lstinline!TYPECASE!} (M2+ runtime typecheck) &
Stub dispatch (compiled) \\
Lua & Metatable lookup + manual C-side checks & Hash lookup in
metatable \\
Guile & Manual \passthrough{\lstinline!SCM!} type predicates & Symbol
lookup in module \\
Tcl & String parse on demand & String comparison for subcommand \\
Python & \passthrough{\lstinline!PyArg\_ParseTuple!} format string & MRO
attribute lookup \\
Kawa & JVM \passthrough{\lstinline!checkcast!} instruction &
\passthrough{\lstinline!invokevirtual!} or reflection \\
JScheme & \passthrough{\lstinline!instanceof!} (runtime) & Reflection +
cached method table \\
CFFI & CLOS \passthrough{\lstinline!typep!} + manual checks & Generic
function dispatch \\
Racket & cpointer tag check & Direct call (no dispatch) \\
\end{longtable}
}

MScheme benefits from M3's runtime type system:
\passthrough{\lstinline!TYPECASE!} is a compiler-supported, type-safe
downcast with O(1) performance (the runtime checks the typecode, which
is an integer). This is essentially the same mechanism as Java's
\passthrough{\lstinline!instanceof!} check, used by Kawa and JScheme on
the JVM.

\subsection{4.3 Automatic Binding
Generation}\label{automatic-binding-generation}

{\def\LTcaptype{none} % do not increment counter
\begin{longtable}[]{@{}
  >{\raggedright\arraybackslash}p{(\linewidth - 6\tabcolsep) * \real{0.0901}}
  >{\raggedright\arraybackslash}p{(\linewidth - 6\tabcolsep) * \real{0.2162}}
  >{\raggedright\arraybackslash}p{(\linewidth - 6\tabcolsep) * \real{0.2613}}
  >{\raggedright\arraybackslash}p{(\linewidth - 6\tabcolsep) * \real{0.4324}}@{}}
\toprule\noalign{}
\begin{minipage}[b]{\linewidth}\raggedright
System
\end{minipage} & \begin{minipage}[b]{\linewidth}\raggedright
Auto-generation?
\end{minipage} & \begin{minipage}[b]{\linewidth}\raggedright
Input
\end{minipage} & \begin{minipage}[b]{\linewidth}\raggedright
Mechanism
\end{minipage} \\
\midrule\noalign{}
\endhead
\bottomrule\noalign{}
\endlastfoot
MScheme & Yes (\passthrough{\lstinline!sstubgen!}) & M3
\passthrough{\lstinline!.i3!} interface files & M3 Toolkit AST parsing +
code generation \\
Tinylisp & Yes (\passthrough{\lstinline!compile\_tli!}) & Hand-written
\passthrough{\lstinline!.tli!} files & Stub compiler generates M2+
stubs \\
Lua & No & -- & Manual C code \\
Guile & Partial (g-wrap) & C headers & Third-party tool, not standard \\
Tcl & Partial (critcl, SWIG) & C headers & Third-party tools \\
Python & Partial (Cython, etc.) & C/C++ headers + annotations &
Third-party tools \\
Kawa & Unnecessary & -- & Shared JVM class loader \\
JScheme & Unnecessary & -- & Reflection on JVM classes at runtime \\
CFFI & No & Manual declarations & Programmer writes
\passthrough{\lstinline!defcfun!} \\
Racket & No & Manual declarations & Programmer writes
\passthrough{\lstinline!define-ffi-definer!} \\
\end{longtable}
}

MScheme's \passthrough{\lstinline!sstubgen!} is notable because it is a
\emph{first-party} tool that ships with the interpreter, processes the
host language's standard interface files, and generates complete
bidirectional bindings. The equivalent in the Python world would be a
tool that reads \passthrough{\lstinline!.pyi!} stub files and generates
complete C extension modules automatically.

The closest comparable systems are SWIG (which processes C/C++ headers
to generate bindings for multiple target languages) and Kawa (which
needs no generator at all because Java and Scheme share the JVM).

\subsection{4.4 GC Coordination}\label{gc-coordination}

{\def\LTcaptype{none} % do not increment counter
\begin{longtable}[]{@{}
  >{\raggedright\arraybackslash}p{(\linewidth - 6\tabcolsep) * \real{0.0917}}
  >{\raggedright\arraybackslash}p{(\linewidth - 6\tabcolsep) * \real{0.1835}}
  >{\raggedright\arraybackslash}p{(\linewidth - 6\tabcolsep) * \real{0.2569}}
  >{\raggedright\arraybackslash}p{(\linewidth - 6\tabcolsep) * \real{0.4679}}@{}}
\toprule\noalign{}
\begin{minipage}[b]{\linewidth}\raggedright
System
\end{minipage} & \begin{minipage}[b]{\linewidth}\raggedright
Host GC
\end{minipage} & \begin{minipage}[b]{\linewidth}\raggedright
Guest GC
\end{minipage} & \begin{minipage}[b]{\linewidth}\raggedright
Coordination
\end{minipage} \\
\midrule\noalign{}
\endhead
\bottomrule\noalign{}
\endlastfoot
MScheme & M3 traced refs & Same GC & None needed -- unified heap \\
Tinylisp & M2+ traced refs & Same GC & None needed -- unified heap \\
Lua & None (C) & Incremental mark-sweep & Explicit: stack + registry
roots \\
Guile & None (C) & Conservative (BDW-GC) & Automatic: scans C stack
conservatively \\
Tcl & None (C) & Refcounting & Manual:
\passthrough{\lstinline!Tcl\_IncrRefCount!}/\passthrough{\lstinline!Tcl\_DecrRefCount!} \\
Python & None (C) & Refcounting + cycle detect & Manual:
\passthrough{\lstinline!Py\_INCREF!}/\passthrough{\lstinline!Py\_DECREF!} \\
Kawa & JVM GC & Same GC & None needed -- unified heap \\
JScheme & JVM GC & Same GC & None needed -- unified heap \\
CFFI & CL GC & Same GC (for CL objects) &
\passthrough{\lstinline!trivial-garbage:finalize!} for foreign ptrs \\
Racket & Precise moving GC & Same GC & Special APIs for C extensions
holding Racket refs \\
\end{longtable}
}

The MScheme/M3, Tinylisp/M2+, Kawa/JVM, and JScheme/JVM systems share a
critical advantage: since both languages use the same garbage collector,
there is no coordination problem. A Scheme closure that captures an M3
object keeps it alive automatically. An M3 data structure containing
Scheme pairs is traced correctly. No reference counting,
prevent-collection guards, or weak-reference bridges.

This eliminates an entire class of bugs that plagues C-based embedding
systems. In CPython, getting reference counts wrong is the single most
common source of C extension bugs. In Lua, forgetting to anchor a value
on the stack or in the registry causes it to be collected while C still
holds a pointer to it.

The fundamental problem runs deeper than API discipline, and
understanding it clarifies why the shared-heap architecture is confined
to a small set of languages: C makes garbage collection nearly
impossible to do correctly. The language permits pointers to be hidden
in ways that are both legal and common in practice, defeating any
collector that tries to find them.

A ``conservative'' collector like BDW-GC works by scanning the C stack,
registers, and static data segments for bit patterns that look like they
might be heap addresses. This is inherently a heuristic, and it fails in
both directions.

\textbf{False negatives (hidden pointers -- the dangerous case).} C
allows a program to store a pointer as an integer
(\passthrough{\lstinline!uintptr\_t!}), XOR two pointers together (XOR
linked lists), offset a pointer by a constant and recover it later,
split a pointer's bytes across non-adjacent struct fields, send it
through a pipe or socket and read it back, or store it in a
memory-mapped file. In all of these cases, the pointer exists but does
not appear on the stack or in a register in a form the collector can
recognize. If the collector runs at that moment, it may conclude the
object is unreachable and free it. The program then recovers the hidden
pointer and dereferences freed memory. This is not a contrived scenario
-- XOR-linked lists save memory in embedded systems, pointer tagging is
used in language runtimes (including CPython's own
\passthrough{\lstinline!ob\_refcnt!} field packing), and integer casts
are routine in serialization and IPC code.

\textbf{False positives (pointer lookalikes -- the slow leak).} On a
32-bit system, roughly 1 in 1000 random integers looks like a valid heap
address. On 64-bit systems the odds are better but not zero, especially
for programs that process file offsets, hash values, or pixel data in
the gigabyte range. Each false positive pins a dead object (and
everything it transitively references) in memory indefinitely. Programs
that run for a long time with large heaps can see unbounded memory
growth from these phantom roots.

\textbf{No solution within C.} The C standard deliberately leaves
pointer representation implementation-defined and permits casts between
pointers and integers. A conforming C program can hide a pointer in a
way that no scanning algorithm can find. This is not a bug in the
language -- it is a consequence of C's design as a portable assembly
language with minimal abstraction over the machine.

Languages like Modula-3 and Java avoid this entirely. In M3, the
compiler emits stack frame descriptors that tell the garbage collector
exactly which stack slots and object fields contain traced references.
The runtime knows the precise layout of every heap object because
allocations go through \passthrough{\lstinline!NEW!}, which records the
type. The collector does not guess: it follows exactly the traced
references and nothing else. No false negatives, false positives, or
heuristics. Java achieves the same property through the JVM's typed
operand stack and GC maps. This is what makes the unified-heap
architecture of MScheme, JScheme, and Kawa possible -- the collector's
precision is guaranteed by the language, not hoped for by convention.

\subsection{4.5 Registration Direction: Push
vs.~Pull}\label{registration-direction-push-vs.-pull}

\textbf{Push} (the host declares what's available to the script):
MScheme, Tinylisp, Lua, Guile, Tcl, Python.

\textbf{Pull} (the script declares what it wants from the host): CFFI,
Racket FFI.

\textbf{Neither} (shared runtime): Kawa, JScheme.

MScheme is firmly push-based, but its push is \emph{automated} via
\passthrough{\lstinline!sstubgen!}. The programmer adds
\passthrough{\lstinline!SchemeStubs("InterfaceName")!} to the makefile,
and the build system generates and compiles all binding code. The effect
feels almost like Kawa's shared runtime: everything visible in M3
interfaces becomes visible in Scheme.

\subsection{4.6 Expressiveness at the
Boundary}\label{expressiveness-at-the-boundary}

MScheme's design allows patterns that are difficult or impossible in
most other systems:

\textbf{M3 objects as first-class Scheme data}. A Scheme list can
contain M3 \passthrough{\lstinline!TextSeq.T!} objects alongside Scheme
pairs and numbers. \passthrough{\lstinline!map!},
\passthrough{\lstinline!filter!}, and \passthrough{\lstinline!for-each!}
work on these mixed lists. The M3 objects participate in Scheme's
structural equality and serialization (via
\passthrough{\lstinline!write!} and pickling).

\textbf{Scheme closures as M3 callbacks}. The
\passthrough{\lstinline!sstubgen!} documentation describes how M3 object
methods can be overridden with Scheme lambdas, enabling a pattern where
M3 objects have individual methods replaced at runtime -- something M3
itself cannot do (methods are fixed at object allocation time).

\textbf{Runtime discovery}. Scheme code can enumerate all registered M3
types, query their supertypes, and list available operations. This
enables Scheme-level metaprogramming over the M3 type system -- writing
generic operations that work on any M3 object based on runtime type
introspection.

\begin{center}\rule{0.5\linewidth}{0.5pt}\end{center}

\section{5. Positioning in the Design
Space}\label{positioning-in-the-design-space}

Plotting the systems on two axes -- \emph{per-call cost} (boundary
overhead) vs.~\emph{setup cost} (how much binding code must be written)
-- reveals three clusters:

\begin{lstlisting}
                    low setup cost
                         |
            Kawa  *      |
                         |
      MScheme  *         |
    Tinylisp  *          |       * JScheme
  -------------------------------------
  low per-call cost      |       high per-call cost
                         |
                         |      * Guile
              * CFFI     |      * Lua
              * Racket   |      * Tcl
                         |      * Python/C
                         |
                    high setup cost
\end{lstlisting}

Kawa achieves the ideal: zero setup, zero per-call cost. But it requires
both languages to run on the JVM.

JScheme shares Kawa's low setup cost (both live on the JVM, no bindings
to write) but pays a significant per-call cost: every Java method call
goes through reflection, while Kawa compiles to direct
\passthrough{\lstinline!invokevirtual!} bytecode.

MScheme and Tinylisp occupy a similar region: both achieve very low
per-call cost through the shared \passthrough{\lstinline!REFANY!} heap,
and both have stub generators that reduce setup cost. MScheme's
\passthrough{\lstinline!sstubgen!} reduces setup cost further by reading
standard \passthrough{\lstinline!.i3!} files directly, while Tinylisp's
\passthrough{\lstinline!compile\_tli!} requires hand-written
\passthrough{\lstinline!.tli!} files. Tinylisp compensates with lower
per-call overhead at the Lisp level (compiled vs. interpreted).

MScheme comes closest to Kawa's position among systems where the host
language has its own native runtime. The shared GC eliminates per-call
cost for value passing, and \passthrough{\lstinline!sstubgen!}
eliminates setup cost for binding generation. The remaining cost is
method dispatch (atom comparison), which is comparable to Lua's
metatable lookup and far cheaper than Python's MRO traversal.

The C-hosted systems (Lua, Guile, Tcl, Python) all pay for the impedance
mismatch between a non-GC host and a GC'd guest. CFFI and Racket
minimize setup cost on the host side (the C library is unmodified) but
pay for it on the scripting side (manual declarations) and at every call
boundary (marshaling).

\begin{center}\rule{0.5\linewidth}{0.5pt}\end{center}

\section{6. Native Code, Not a Virtual
Machine}\label{native-code-not-a-virtual-machine}

The preceding analysis groups MScheme with JScheme and Kawa as ``shared
heap'' systems. But there is a further distinction that separates
MScheme from its JVM counterparts, and connects it to a much smaller
lineage.

\subsection{6.1 Why This Distinction
Matters}\label{why-this-distinction-matters}

JScheme and Kawa achieve glueless interoperability on the JVM. But the
JVM is a virtual machine: Java source compiles to bytecode, which is
then interpreted or JIT-compiled by the VM at runtime. The GC, runtime
type information, and universal reference type
(\passthrough{\lstinline!java.lang.Object!}) are services provided by
the VM, not by the compiled program itself.

Modula-3 is different. The cm3 compiler emits native machine code --
x86, x86-64, ARM64, SPARC, PA-RISC, PowerPC -- directly, with no
intermediate bytecode and no virtual machine. The generated code runs as
ordinary machine instructions. Yet Modula-3 still provides precise
garbage collection, runtime type discrimination
(\passthrough{\lstinline!TYPECASE!}, \passthrough{\lstinline!ISTYPE!},
\passthrough{\lstinline!TYPECODE!}), and a universal traced reference
type (\passthrough{\lstinline!REFANY!}). It achieves this by having the
compiler emit metadata alongside the native code: stack frame
descriptors that tell the runtime exactly which stack slots contain
traced references, and type descriptors that record the layout of every
heap-allocated type.

This means MScheme's Scheme interpreter runs as native machine code
calling into native machine code, with no VM layer in between. When
Scheme code calls an M3 method, it is a direct procedure call -- the
same calling convention as any M3-to-M3 call. When the garbage collector
runs, it traces M3 objects and Scheme pairs with the same algorithm,
using the same compiler-generated stack maps. There is no bytecode
interpreter, JIT warmup, or tiered compilation.

\subsection{6.2 How Rare Is This
Architecture?}\label{how-rare-is-this-architecture}

The combination of properties that makes MScheme possible is unusual:
(1) ahead-of-time compilation to native code, (2) precise tracing GC,
(3) runtime type information with a universal reference type, and (4) an
embedded interpreter for a \emph{different} language that shares the
host's values and GC directly, with no marshaling.

Very few languages in history have provided (1)-(3) simultaneously. The
lineage leading to Modula-3 -- Mesa → Cedar → Modula-2+ → Modula-3 -- is
one of the few that does (see section 8.1).

A survey of other natively-compiled GC'd languages shows why MScheme's
architecture is hard to replicate:

\textbf{Common Lisp (SBCL, CMUCL, etc.).} The strongest analogue. SBCL
compiles to native x86-64 and ARM64 code and has a precise generational
GC. Its interpreter (when enabled) shares the exact same heap, the same
tagged-pointer value representation, and the same collector as compiled
code. A Scheme interpreter written in Common Lisp would be
architecturally identical to MScheme. However, CL's interpreter is
self-hosting -- it interprets Common Lisp, not a different language. The
system is designed for mixed-mode execution of a single language, not
for embedding a scripting language in a systems language.

\textbf{Smalltalk (StrongTalk, Cog/Spur).} JIT-compiles methods to
native code; the interpreter and compiled code share a single object
space with no marshaling boundary. Again self-hosting: the interpreted
and compiled languages are the same Smalltalk. Compilation is JIT (at
runtime), not AOT (ahead of time).

\textbf{Eiffel (EiffelStudio's Melting Ice).} Compiles Eiffel to C and
thence to native code (``frozen'' code), but also maintains bytecode-
interpreted ``melted'' code for rapid development. Frozen and melted
code share the same heap and objects. Self-hosting: the interpreted
language is Eiffel.

\textbf{Julia.} JIT-compiles to native code via LLVM;
\passthrough{\lstinline!eval!} operates in the same runtime with the
same values. Self-hosting; JIT, not AOT.

\textbf{Go (Yaegi).} Compiles to native code and has precise GC, but
Yaegi (a Go interpreter written in Go) wraps values in
\passthrough{\lstinline!reflect.Value!} at every boundary -- a form of
marshaling. Go also lacks a universal reference type;
\passthrough{\lstinline!interface\{\}!} is a two-word pair, not a traced
pointer.

\textbf{Oberon.} Compiles to native code and has GC, but lacks a
universal reference type and has no embedded interpreter -- it uses
dynamic loading of compiled modules instead.

\textbf{D, Dylan, OCaml.} All compile to native code and have GC, but
none has an embedded interpreter for a different language that shares
the host heap without marshaling. OCaml's bytecode toplevel shares the
heap with loaded \passthrough{\lstinline!.cmo!} modules, but the
native-code compiler path has no analogous mechanism for a different
language.

The pattern that emerges is that self-hosting interpreters in
natively-compiled languages are not uncommon (CL, Smalltalk, Eiffel,
Julia all have them), but an interpreter for a \emph{different} language
that shares the host's heap with zero marshaling overhead is rare. The
obstacle is not theoretical -- any language with GC, RTTI, and a
universal reference type could support it -- but practical: very few
such languages exist, and in those that do, the idea has seldom been
pursued. The one important exception -- Tinylisp (section 3.9) --
exploits the same mechanism in Modula-2+, the direct ancestor of
Modula-3.

\begin{center}\rule{0.5\linewidth}{0.5pt}\end{center}

\section{7. Design Philosophy: Value Types and the Lisp
Machine}\label{design-philosophy-value-types-and-the-lisp-machine}

The SSTUBGEN.TXT design document explicitly positions the system's goal
as making Scheme a viable ``main program'' language for large M3
applications: the libraries are in M3, the orchestration is in Scheme.
This is the same pattern later popularized by Lua in game engines and
Python in scientific computing -- but MScheme achieves it with
significantly less impedance mismatch, thanks to the shared GC and type
system.

Why do so few languages provide the features that make this architecture
possible? The answer leads to a deeper question about language design
and the relationship between MScheme's architecture and the Lisp
Machine.

The Lisp Machine is the ultimate precedent for glueless
interoperability: Lisp \emph{was} the systems language and no
interoperability boundary existed at all. But the Lisp Machine also
illustrates the cost of the approach the Algol family avoids. As Nelson
observes in \emph{Systems Programming with Modula-3} (1991):

\begin{quote}
At the other extreme are languages like Lisp, ML, Smalltalk, and CLU,
whose programming models originate from mathematics. {[}\ldots{]} These
languages have beautiful programming models, but they tend to be
difficult to implement efficiently, because the uniform treatment of
values in the programming model invites a runtime system in which values
are uniformly represented by pointers. If the implementer doesn't take
steps to avoid it, as simple a statement as
\passthrough{\lstinline!n := n + 1!} could require an allocation, a
method lookup, or both.
\end{quote}

The Lisp Machine solved this in hardware. The Symbolics 3600 (Moon, ISCA
1985) and its successor the Ivory (I-Machine) use tagged words: every
machine word carries a data type tag (4 bits on the 3600, 6 bits on the
Ivory) alongside 32 bits of data or address. Twenty-two of the Ivory's
type codes designate \emph{immediate} values -- fixnums,
single-precision floats, characters, small ratios -- where the data is
encoded directly in the word. Fixnum arithmetic is a
register-to-register operation; \passthrough{\lstinline!n + 1!} requires
no allocation and no GC involvement. The remaining 33 type codes
designate \emph{pointer} values -- lists, arrays, symbols, closures,
bignums -- where the 32-bit field is a heap address. CDR-coding
compresses list cells to roughly half the space of naive linked lists.
Hardware-assisted ephemeral garbage collection reduces the cost of
short-lived heap allocations. Common Lisp's
\passthrough{\lstinline!dynamic-extent!} declaration allows the compiler
to stack-allocate objects known not to escape their enclosing scope.
With all of these mechanisms, Lisp Machines do handle fixnum-intensive
code efficiently -- but they achieve it through \emph{dedicated
hardware} that conventional machines lack.

Tagged-pointer hardware has a longer history -- the Burroughs B5000
(1961) and B6500/B6700 (1969) tagged every word for machine-level
integrity -- but the Lisp Machines were the first to tag by
\emph{user-level data type} (fixnum, cons, symbol, closure), enabling
the hardware to dispatch on language-level types directly.

The Algol family, to which Modula-3 belongs, takes the opposite
approach: instead of solving the problem in hardware, the \emph{language
itself} distinguishes value types from reference types. In Modula-3,
\passthrough{\lstinline!INTEGER!}, \passthrough{\lstinline!REAL!},
\passthrough{\lstinline!BOOLEAN!}, \passthrough{\lstinline!CHAR!},
records, and fixed arrays are value types -- stack-allocated, unboxed,
and invisible to the garbage collector. Only reference types
(\passthrough{\lstinline!REF!}, \passthrough{\lstinline!OBJECT!},
\passthrough{\lstinline!TEXT!}) participate in the traced heap. This
means that most of a typical M3 systems program -- local variables, loop
counters, parameter passing, record manipulation -- never touches the GC
at all. The garbage collector runs only for heap-allocated reference
types, and the compiler needs no special optimization or declarations to
achieve this: the type system guarantees it statically.

This is exactly the property that makes MScheme's architecture practical
for systems programming. The Scheme interpreter and its data structures
(pairs, closures, environments) live in the traced heap and are managed
by the GC, but the M3 code that the Scheme stubs call into can be
dominated by stack-allocated value types. The boundary between ``GC'd
scripting world'' and ``stack-allocated systems world'' is not a
marshaling barrier -- it is the ordinary Modula-3 type system, which
\passthrough{\lstinline!sstubgen!} navigates automatically by generating
\passthrough{\lstinline!ToScheme!}/\passthrough{\lstinline!ToModula!}
converters that box and unbox value types at the boundary while passing
reference types through unchanged. (JScheme and Kawa on the JVM face a
version of the same issue -- the JVM boxes primitives like
\passthrough{\lstinline!int!} to \passthrough{\lstinline!Integer!} --
but Java's autoboxing covers only a fixed set of primitive types, while
\passthrough{\lstinline!sstubgen!} handles the full M3 type system
including user-defined records, enumerations, subranges, and
multidimensional arrays.)

The Lisp Machine achieved glueless interoperability by making Lisp the
only language. JScheme approximated it for Java on the JVM. MScheme
approximates it for Modula-3 on conventional hardware -- one type alias
(\passthrough{\lstinline!T = REFANY!}), a language that distinguishes
values from references, and a code generator that handles the rest.

\begin{center}\rule{0.5\linewidth}{0.5pt}\end{center}

\section{8. Historical Lineage}\label{historical-lineage}

\subsection{8.1 The Mesa/Cedar/Modula-2+/Modula-3 Language
Family}\label{the-mesacedarmodula-2modula-3-language-family}

The language features that make MScheme and Tinylisp possible --
\passthrough{\lstinline!REFANY!}, precise tracing GC,
\passthrough{\lstinline!TYPECASE!} -- trace a specific lineage through
four decades of systems programming language design.

\textbf{Mesa} (Xerox PARC, 1970s) was a systems language for the Alto
and later the Dorado. It compiled to native code and had strong typing,
but no garbage collection and no universal reference type.

\textbf{Cedar} (Xerox PARC, early 1980s) extended Mesa with garbage
collection, \passthrough{\lstinline!REF ANY!}, and runtime type
information. Paul Rovner's 1985 technical report documents the addition
of these features. Cedar had the infrastructure to support an
MScheme-like system -- a universal traced reference type and a precise
GC -- but no one built a full scripting interpreter on top of it. The
debugger could evaluate expressions in a stopped frame, but that is not
the same as a general-purpose embedded language.

\textbf{Modula-2+} (DEC Systems Research Center, mid-1980s) was DEC
SRC's dialect of Modula-2, adding garbage collection, exception
handling, opaque types, \passthrough{\lstinline!REFANY!}, and
\passthrough{\lstinline!TYPECASE!} -- the features Cedar had pioneered,
adapted to a Wirth-family language. It was the implementation language
for SRC's systems infrastructure and the language in which Tinylisp was
written.

\textbf{Modula-3} (DEC SRC / Olivetti Research Center, 1989) refined
Modula-2+ into a clean design with a formal specification. The Modula-3
Report (Cardelli, Donahue, Glassman, Jordan, Kalsow, and Nelson, 1989)
codified the type system, the \passthrough{\lstinline!REFANY!}
mechanism, and the \passthrough{\lstinline!TYPECASE!} dispatch that make
MScheme's architecture possible.

\subsection{8.2 From JScheme to MScheme}\label{from-jscheme-to-mscheme}

Peter Norvig published JScheme in 1998 as a compact demonstration that a
useful Scheme interpreter could be written in \textasciitilde1700 lines
of Java. The key architectural insight --
\passthrough{\lstinline!java.lang.Object!} as the universal Scheme value
-- was implicit in the choice of Java as the implementation language.
Ken Anderson (BBN) and Tim Hickey (Brandeis) later expanded JScheme with
dot-notation syntax, overload resolution, and dynamic proxy support for
Java interfaces.

MScheme was written in 2007-2008 as a line-for-line translation of
JScheme into Modula-3. The translation preserved JScheme's fundamental
principle (\passthrough{\lstinline!Object!} -\textgreater{}
\passthrough{\lstinline!REFANY!}) but added two things JScheme lacks: an
automatic stub generator (\passthrough{\lstinline!sstubgen!}) and a
runtime type introspection API. MScheme has been used in financial
trading, scientific computing, and hardware design.

The MScheme design predates many of the modern FFI systems (Racket's FFI
was redesigned around 2010; Python's stable ABI came in 3.2/2011;
Guile's foreign object types replaced SMOBs in 2.0/2011).

MScheme was not aware of Tinylisp; the parallel development reflects the
fact that the architecture is essentially forced once the host language
provides \passthrough{\lstinline!REFANY!} and GC.

\subsection{8.3 The sstubgen and m3tk
Provenance}\label{the-sstubgen-and-m3tk-provenance}

\passthrough{\lstinline!sstubgen!} derives directly from the Network
Objects stub generator (\passthrough{\lstinline!m3-comm/stubgen!})
designed by Susan Owicki and Ted Wobber for Birrell et al.'s distributed
object system (SRC Report 115, 1994). It is not merely modeled on the
network objects \passthrough{\lstinline!stubgen!} -- it shares code.
Several source files are identical between the two packages
(\passthrough{\lstinline!FRefRefTbl!},
\passthrough{\lstinline!TypeNames!},
\passthrough{\lstinline!ValueProc!}), and others
(\passthrough{\lstinline!AstToType!}, \passthrough{\lstinline!Type!},
\passthrough{\lstinline!Value!}, \passthrough{\lstinline!StubGenTool!},
\passthrough{\lstinline!StubUtils!}) are extended copies.

M3tk is a comprehensive Modula-3 front-end toolkit providing a generic
AST framework, parser, semantic analyzer, and support for building
compiler extensions. It was written by Mick Jordan, originally at the
Olivetti Research Center (ORC) in Menlo Park, California, and later at
DEC SRC, and described in his 1990 paper ``An Extensible Programming
Environment for Modula-3'' (ACM SDE 4). M3tk began life as part of
Olivetti's full Modula-3 compiler; the cm3 compiler (Bill Kalsow's SRC
compiler lineage) does not use m3tk itself, but the Network Objects
\passthrough{\lstinline!stubgen!} does, and
\passthrough{\lstinline!sstubgen!} inherits that dependency. M3tk is
organized into sub-packages (\passthrough{\lstinline!m3tk-ast!},
\passthrough{\lstinline!m3tk-syn!}, \passthrough{\lstinline!m3tk-sem!},
\passthrough{\lstinline!m3tk-fe!}, etc.) and is used by both
\passthrough{\lstinline!stubgen!} and
\passthrough{\lstinline!sstubgen!}, as well as by the language server
and other M3 development tools.

\begin{center}\rule{0.5\linewidth}{0.5pt}\end{center}

\section{Acknowledgments}\label{acknowledgments}

Paul McJones provided detailed comments on an earlier draft, including
the reference to Ellis's Tinylisp and corrections to the discussion of
m3tk, stub naming conventions, and the role of the
\passthrough{\lstinline!interp!} and \passthrough{\lstinline!excH!}
parameters in generated code.

\begin{center}\rule{0.5\linewidth}{0.5pt}\end{center}

\section{Bibliography}\label{bibliography}

\subsection{Modula-3}\label{modula-3}

\begin{itemize}
\item
  Cardelli, L., Donahue, J., Glassman, L., Jordan, M., Kalsow, B., and
  Nelson, G. ``Modula-3 Report (revised).'' DEC Systems Research Center,
  Report 52, November 1989.
\item
  Cardelli, L., Donahue, J., Jordan, M., Kalsow, B., and Nelson, G.
  ``The Modula-3 Type System.'' In \emph{Conference Record of the
  Sixteenth Annual ACM Symposium on Principles of Programming
  Languages}, pp.~202--212, 1989.
\item
  Cardelli, L. ``Typeful Programming.'' DEC Systems Research Center,
  Report 45, May 1989. Revised version in E.J. Neuhold and M. Paul
  (eds.), \emph{Formal Description of Programming Concepts},
  Springer-Verlag, 1991.
\item
  Nelson, G. (ed.) \emph{Systems Programming with Modula-3.} Prentice
  Hall, 1991.
\item
  Harbison, S.P. \emph{Modula-3.} Prentice Hall, 1992.
\item
  Modula3/CM3: Critical Mass Modula-3. https://github.com/modula3/cm3
\end{itemize}

\subsection{Network Objects and stub
generation}\label{network-objects-and-stub-generation}

\begin{itemize}
\item
  Birrell, A.D., Nelson, G., Owicki, S., and Wobber, E. ``Network
  Objects.'' DEC Systems Research Center, Report 115, February 1994.
  Also in \emph{Proceedings of the 14th ACM Symposium on Operating
  Systems Principles} (SOSP), December 1993; and
  \emph{Software--Practice and Experience}, 25(S4):87--130, December
  1995.
\item
  Birrell, A.D., Evers, D., Nelson, G., Owicki, S., and Wobber, E.
  ``Distributed Garbage Collection for Network Objects.'' DEC Systems
  Research Center, Report 116, December 1993.
\end{itemize}

\subsection{M3 Toolkit (m3tk)}\label{m3-toolkit-m3tk}

\begin{itemize}
\tightlist
\item
  Jordan, M. ``An Extensible Programming Environment for Modula-3.'' In
  \emph{Proceedings of the Fourth ACM SIGSOFT Symposium on Software
  Development Environments} (SDE 4), Irvine, California, December

  \begin{enumerate}
  \def\labelenumi{\arabic{enumi}.}
  \setcounter{enumi}{1989}
  \tightlist
  \item
    \emph{ACM SIGSOFT Software Engineering Notes}, 15(6):66--76.
  \end{enumerate}
\item
  Jordan, M. and Robinson, P. ``A Programming Environment for
  Modula-2.'' \emph{Software Engineering Journal}, 3(4):119--126, July

  \begin{enumerate}
  \def\labelenumi{\arabic{enumi}.}
  \setcounter{enumi}{1987}
  \tightlist
  \item
  \end{enumerate}
\end{itemize}

\subsection{\texorpdfstring{MScheme, \texttt{sstubgen}, and
\texttt{cspc}}{MScheme, sstubgen, and cspc}}\label{mscheme-sstubgen-and-cspc}

\begin{itemize}
\item
  MScheme source code: \passthrough{\lstinline!cm3/m3-scheme/mscheme/!}.
  In the CM3 repository.
\item
  SSTUBGEN.TXT: Design document for the Scheme stub generator.
  \passthrough{\lstinline!cm3/m3-scheme/sstubgen/src/SSTUBGEN.TXT!}.
\item
  SchemeM3 README: Embedding guide.
  \passthrough{\lstinline!cm3/m3-scheme/modula3scheme/src/README!}.
\item
  \passthrough{\lstinline!cspc!}: CSP compiler for asynchronous VLSI
  design. A Modula-3 program embedding MScheme; the compiler itself is
  written in Scheme, using \passthrough{\lstinline!sstubgen!}-generated
  stubs to access M3 libraries.
  \passthrough{\lstinline!intel-async/async-toolkit/m3utils/m3utils/csp/src/!}.
  In the async-toolkit repository.
\end{itemize}

\subsection{JScheme}\label{jscheme}

\begin{itemize}
\item
  Norvig, P. ``JScheme: Scheme Implemented in Java.''
  https://norvig.com/jscheme.html, 1998.
\item
  Norvig, P. ``JScheme: Design Decisions.''
  https://norvig.com/jscheme-design.html
\item
  Anderson, K.R., Hickey, T.J., and Norvig, P. ``JScheme.''
  https://www.cs.brandeis.edu/\textasciitilde tim/jscheme/main.html.
  Brandeis University / BBN Technologies, 2003.
\end{itemize}

\subsection{Tinylisp and Modula-2+}\label{tinylisp-and-modula-2}

\begin{itemize}
\item
  Ellis, J. ``Tinylisp Reference Manual.'' Systems Research Center,
  Digital Equipment Corporation, January 26, 1989.
  https://softwarepreservation.computerhistory.org/LISP/dec/Ellis-Tinylisp-1989.pdf
\item
  Modula-2+ manuals and related materials.
  https://softwarepreservation.computerhistory.org/modula2+/
\end{itemize}

\subsection{Cedar and Modula-3
ancestry}\label{cedar-and-modula-3-ancestry}

\begin{itemize}
\item
  Rovner, P. ``On Adding Garbage Collection and Runtime Types to a
  Strongly-Typed, Statically-Checked, Concurrent Language.'' Xerox PARC,
  CSL-84-7, July 1985.
\item
  Lampson, B.W. ``A Description of the Cedar Language.'' Xerox PARC,
  CSL-83-15, December 1983.
\end{itemize}

\subsection{Lua}\label{lua}

\begin{itemize}
\item
  Ierusalimschy, R., de Figueiredo, L.H., and Celes, W. \emph{Lua 5.4
  Reference Manual.} https://www.lua.org/manual/5.4/
\item
  Ierusalimschy, R. \emph{Programming in Lua}, 4th edition. Lua.org,

  \begin{enumerate}
  \def\labelenumi{\arabic{enumi}.}
  \setcounter{enumi}{2015}
  \tightlist
  \item
  \end{enumerate}
\item
  Ierusalimschy, R., de Figueiredo, L.H., and Celes, W. ``The Evolution
  of Lua.'' In \emph{Proceedings of the Third ACM SIGPLAN Conference on
  History of Programming Languages} (HOPL III), 2007.
\end{itemize}

\subsection{GNU Guile}\label{gnu-guile}

\begin{itemize}
\tightlist
\item
  \emph{GNU Guile Reference Manual}, version 3.0.
  https://www.gnu.org/software/guile/manual/
\end{itemize}

\subsection{Tcl/Tk}\label{tcltk}

\begin{itemize}
\item
  Ousterhout, J.K. \emph{Tcl and the Tk Toolkit.} Addison-Wesley, 1994.
\item
  \emph{Tcl/Tk Documentation.} https://www.tcl-lang.org/doc/
\end{itemize}

\subsection{CPython}\label{cpython}

\begin{itemize}
\item
  \emph{Python/C API Reference Manual.} https://docs.python.org/3/c-api/
\item
  \emph{Extending and Embedding the Python Interpreter.}
  https://docs.python.org/3/extending/
\end{itemize}

\subsection{Kawa Scheme}\label{kawa-scheme}

\begin{itemize}
\item
  Bothner, P. ``Kawa -- Compiling Dynamic Languages to the Java VM.'' In
  \emph{Proceedings of the USENIX Annual Technical Conference}, FREENIX
  Track, 1998.
\item
  \emph{The Kawa Scheme Language.} https://www.gnu.org/software/kawa/
\end{itemize}

\subsection{Common Lisp CFFI}\label{common-lisp-cffi-1}

\begin{itemize}
\item
  \emph{CFFI: The Common Foreign Function Interface.}
  https://cffi.common-lisp.dev/
\item
  \emph{CFFI User Manual.}
  https://cffi.common-lisp.dev/manual/cffi-manual.html
\end{itemize}

\subsection{Racket FFI}\label{racket-ffi-1}

\begin{itemize}
\item
  Barzilay, E. and Orlovsky, D. ``Foreign Interface for PLT Scheme.'' In
  \emph{Proceedings of the Fifth ACM SIGPLAN Workshop on Scheme and
  Functional Programming}, 2004.
\item
  \emph{The Racket Foreign Interface.}
  https://docs.racket-lang.org/foreign/
\end{itemize}

\subsection{Conservative garbage
collection}\label{conservative-garbage-collection}

\begin{itemize}
\item
  Boehm, H.-J. and Weiser, M. ``Garbage Collection in an Uncooperative
  Environment.'' \emph{Software--Practice and Experience},
  18(9):807--820, September 1988.
\item
  Boehm, H.-J., Demers, A.J., and Shenker, S. ``Mostly Parallel Garbage
  Collection.'' In \emph{Proceedings of the ACM SIGPLAN '91 Conference
  on Programming Language Design and Implementation}, pp.~157--164,
  1991.
\end{itemize}

\subsection{SBCL}\label{sbcl}

\begin{itemize}
\item
  \emph{SBCL User Manual.} https://www.sbcl.org/manual/
\item
  Steel Bank Common Lisp source. https://github.com/sbcl/sbcl
\end{itemize}

\subsection{Eiffel}\label{eiffel}

\begin{itemize}
\item
  Meyer, B. \emph{Object-Oriented Software Construction}, 2nd edition.
  Prentice Hall, 1997.
\item
  ``Melting Ice Technology.''
  https://www.eiffel.org/doc/eiffelstudio/Melting\_Ice\_Technology
\end{itemize}

\subsection{Smalltalk and StrongTalk}\label{smalltalk-and-strongtalk}

\begin{itemize}
\item
  Bracha, G. and Griswold, D. ``Strongtalk: Typechecking Smalltalk in a
  Production Environment.'' In \emph{Proceedings of the ACM Conference
  on Object-Oriented Programming, Systems, Languages and Applications}
  (OOPSLA), pp.~215--230, 1993.
\item
  Ungar, D. and Smith, R.B. ``Self: The Power of Simplicity.'' In
  \emph{Proceedings of OOPSLA '87}, pp.~227--241, 1987.
\end{itemize}

\subsection{Oberon}\label{oberon}

\begin{itemize}
\tightlist
\item
  Wirth, N. and Gutknecht, J. \emph{Project Oberon: The Design of an
  Operating System and Compiler.} ACM Press / Addison-Wesley, 1992.
  Revised edition (\emph{The Design of an Operating System, a Compiler,
  and a Computer}), 2013.
\end{itemize}

\subsection{Lisp Machines and tagged
architectures}\label{lisp-machines-and-tagged-architectures}

\begin{itemize}
\item
  Moon, D.A. ``Architecture of the Symbolics 3600.'' In
  \emph{Proceedings of the 12th Annual International Symposium on
  Computer Architecture} (ISCA), pp.~76--83, 1985.
\item
  Organick, E.I. \emph{Computer System Organization: The B5700/B6700
  Series.} Academic Press, 1973.
\item
  Levy, H.M. \emph{Capability-Based Computer Systems.} Digital Press,

  \begin{enumerate}
  \def\labelenumi{\arabic{enumi}.}
  \setcounter{enumi}{1983}
  \tightlist
  \item
    Chapter 2 covers the Burroughs B5000 and B6500 tagged architectures.
  \end{enumerate}
\end{itemize}

\subsection{Other systems cited}\label{other-systems-cited}

\begin{itemize}
\item
  Traefik. ``Yaegi: Another Elegant Go Interpreter.''
  https://github.com/traefik/yaegi
\item
  ``Hint: Runtime Haskell Interpreter.''
  https://hackage.haskell.org/package/hint
\item
  Flatt, M., Findler, R.B., and PLT. \emph{The Racket Guide.}
  https://docs.racket-lang.org/guide/
\end{itemize}

\end{document}
